% concatenated from main.tex, preamble.tex, section_files/1_Introduction.tex, section_files/2_General_GE_discussion.tex, section_files/3_GWM_and_RWM_in_1D_R.tex, section_files/5_Discussion.tex, section_files/6_appendix_concise.tex
\documentclass{article}
\usepackage{amsmath,amssymb,graphicx,setspace,amsthm,float,url}
\usepackage{subcaption}
\usepackage{mathtools,blindtext, geometry}
\usepackage{tabularx}
\usepackage{xcolor}
\usepackage[hidelinks]{hyperref}
\usepackage{algorithm}
\usepackage{algpseudocode}
\usepackage{tikz}  % for graph generated
\usepackage[affil-it]{authblk}
\usepackage[numbers,square,sort&compress]{natbib}


% theorem stuff

%% This style keeps the body text upright
\theoremstyle{definition}
\newtheorem{definition}{Definition}[section]
\newtheorem{example}[definition]{Example}
\newtheorem{remark}[definition]{Remark}


%% This style italicizes the body text
\theoremstyle{plain}
\newtheorem{assumption}{Assumption}[section]
\newtheorem{theorem}[definition]{Theorem}
\newtheorem{proposition}[definition]{Proposition}
\newtheorem{lemma}[definition]{Lemma}
\newtheorem{corollary}[definition]{Corollary}
\newtheorem{conjecture}[definition]{Conjecture}
\newtheorem{statement}[definition]{Statement}



\newenvironment{solution}{\renewcommand\qedsymbol{$\blacksquare$}\begin{proof}[Solution]}
{\end{proof}}

% Comments command
\newcommand{\sjlcom}[1]{\textcolor{teal}{#1}}


\title{\huge A note on diffusive/random-walk behaviour in Metropolis--Hastings algorithms}

% 1. Use manual symbols in the author names
\author{Yuxin Liu$^{*, \dagger}$}
\author{Peiyi Zhou$^{*, \ddagger}$}
\author{Samuel Livingstone$^{\S}$}

\affil{Department of Statistical Science, University College London}

\begin{document}

\maketitle

\insert\footins{\noindent\footnotesize $^*$Equal contribution.}
\insert\footins{\noindent\footnotesize $^\dagger$\texttt{y.liu.24@ucl.ac.uk}; Corresponding author}
\insert\footins{\noindent\footnotesize $^\ddagger$\texttt{peiyi.zhou.25@ucl.ac.uk}; Corresponding author}
\insert\footins{\noindent\footnotesize $^\S$\texttt{samuel.livingstone@ucl.ac.uk}}

\begin{abstract}
We prove a general result that if a Metropolis--Hastings algorithm has a proposal that is not geometrically ergodic and the acceptance rate approaches unity at a suitable rate as the state variable becomes large, then the Metropolised chain will also not be geometrically ergodic. Our conditions seem stronger than might be expected, but are shown to be necessary through a counterexample. We then turn our attention to the random walk and guided walk Metropolis algorithms. We show that if the target distribution has polynomial tails the latter converges at twice the polynomial rate of the former, but that if instead the target distribution has strictly convex potential then the random walk Metropolis behaves as a $1/2$-lazy version of the guided walk Metropolis when the state variable is large, and therefore moves at a similar (ballistic) speed.
\end{abstract}

%\tableofcontents

%=========================================
% Create your .tex file in Notes by copy-pasting your main.tex
% then simply \input{Notes/your main.tex}
% if you have references, simply add it into refs.bib
\section{Introduction}

Sampling algorithms and Monte Carlo methods based on Markov processes are an indispensable tool in modern computational research.  A prominent example is the Metropolis--Hastings algorithm, a meta-algorithm from which many popular approaches to Markov chain Monte Carlo can be derived, either as a particular example or in some appropriate limit.  The most basic form of Metropolis--Hastings algorithm will produce a $\pi$-reversible Markov chain, where $\pi$ is the target distribution from which samples are desired, but non-reversible algorithms can straightforwardly be constructed by combining Metropolis--Hastings steps with others in a cyclical manner, leading to a rich suite of methods that has captured the attention of both researchers and practitioners for many decades.

A key question of interest is how to ensure that Markov chain algorithms mix quickly, so that samples and ergodic averages are representative of the target distribution $\pi$.  
Chains that are said to exhibit `random walk' or `diffusive' behaviour take small, directionless steps and uncover little about $\pi$. Such chains are generally shunned because they often mix slowly.  A common approach to reducing random walk behaviour is to augment the state variable with a momentum, typically creating a non-reversible algorithm.  It is known that ergodic averages constructed from a non-reversible Markov chain $P$ will have at least as small an asymptotic variance as a reversible chain formed as $(P + P^*)/2$, where $P^*$ is the $L^2(\pi)$-adjoint of $P$, but the degree of improvement is dependent on both $P$ and $\pi$ \cite{andrieu2016random}.  Mixing times can also be reduced by introducing momentum, but also increased arbitrarily if care is not taken \cite{diaconis2000analysis}.

Diffusive behaviour more generally is an extensively studied phenomenon \citep{metzler2000random}.  A simple categorisation for a stochastic process $(X_t)_{t \geq 0}$ on $\mathbb{R}^d$ is to consider the displacement from its initial value $X_0 := x$ for suitably small $t$.  Typically 
\[
\mathbb{E}\|X_t - x\| \sim t^\alpha,
\]
for some $\alpha > 0$.  The case $\alpha = 1/2$ indicates \textit{diffusive} behaviour, which is characteristic of Brownian motion. If $\alpha < 1/2$ then the process is called \textit{subdiffusive}, and if $\alpha > 1/2$ it is \textit{superdiffusive}.  In the special case $\alpha = 1$ it is called \textit{ballistic}. In any instance that $\alpha \neq 1/2$ the process is said to exhibit \textit{anomalous diffusion}, which is an intense area of mathematical study \citep{metzler2000random, bouchaud1990anomalous}.

A reversible Markov chain Monte Carlo sampling algorithm can behave in a diffusive manner if the transition kernel places most of its mass in a small neighbourhood $N(x)$ of the current state $x$, and the target distribution formally satisfies $\pi(dx) \approx \pi(dy)$ for any $y \in N(x)$. In this case the detailed balance equations dictate that $P(x,dy) \approx P(y,dx)$, meaning the chain will behave in a diffusive manner when this property is combined with the locality of the transition.  A non-reversible algorithm no longer satisfies the detailed balance equations, and so even a transition kernel that still concentrates on $N(x)$ can induce anomalous, often ballistic, behaviour in the same setting.  Reversible Markov chains need not always behave diffusively, however, as the transition is not always restricted to a small neighbourhood of the current state, and the target distribution is not always `flat' in the region that a typical transition is made.  In this case the benefits of non-reversibility are less clear.

We highlight in this article cases in which reversible Metropolis--Hastings algorithms can exhibit diffusive behaviour, and others in which they behave just as non-reversible alternatives do, the key differentiator being the form of $\pi$.  Recall that a single step of a Metropolis--Hastings algorithm consists of drawing a candidate next state for the Markov chain $Y \sim Q(X_i,\cdot)$, where $Q$ is some kernel, and setting $X_{i+1} = Y$ with probability $\alpha(X_i,Y)$, where $\alpha$ is chosen to enforce $\pi$-reversibility.  We begin in Section \ref{sec:MH_high_accept} with a very general result showing that if typical candidate next steps result in $\alpha(X_i,Y) \approx 1$ when $\|X_i\|$ is large, and $Q$ itself is not geometrically ergodic, then the Metropolis--Hastings transition $P$ will typically not be either.  This is a common scenario in which diffusive behaviour can often be observed, because $Q$ alone does not promote fast mixing and $\alpha$ does not influence the dynamics to the degree that $P$ behaves any differently to $Q$.  Surprisingly, however, the condition that $\lim_{\|x\|\to\infty}\int \alpha(x,y)Q(x,dy) = 1$ alone is not sufficient for such a result; we provide a counterexample in which it holds and $Q$ is not geometrically ergodic but $P$ in fact is.  We also discuss extensions in which $\alpha \not\to 1$ but nonetheless the distribution becomes asymptotically flat in some directions, and connect this to the concept of acceleration in sampling and optimisation \cite{oberdorster2025accelerated, bubeck2015convex}.  

In Section \ref{sec: RWM and GWM on 1D R} we focus on two specific algorithms, the random walk and guided walk Metropolis, which can be viewed as a canonical reversible Metropolis--Hastings algorithm and its natural non-reversible counterpart.  If $\pi$ has a density with polynomial tails then it is well-known that the random walk Metropolis can be diffusive while the guided walk exhibits ballistic motion.  We prove that this difference is consequential, in that the guided walk on $\mathbb{R}$ exhibits a faster polynomial rate of convergence to equilibrium.  We then consider the case in which $\pi$ has a strictly log-concave density, meaning that the tails are lighter. In this case the behaviour differs markedly; in the limit $|x|\to\infty$, the random walk Metropolis behaves as a $1/2$-lazy version of the guided walk Metropolis.  Despite $Q$ being a random walk, the acceptance rate qualitatively changes the transition resulting in very little difference in the behaviour of the reversible and non-reversible algorithms during the transient phase, with both exhibiting ballistic motion.


\section{Metropolis--Hastings \& the high-acceptance regime}
\label{sec:MH_high_accept}

The Metropolis--Hastings algorithm provides a simple recipe to construct a Markov transition kernel. Given the current state $X_i$ a candidate kernel $Q$ is used to generate $Y \sim Q(X_i,\cdot)$.  After this, the next state $X_{i+1}$ is set to be $Y$ with probability $\alpha(X_i,Y)$.  There are infinitely many possible choices for $\alpha$, but the most common when both $\pi$ and $Q$ possess densities $\pi(x)$ and $q(x,y)$ with respect to a common dominating measure is $\alpha(x,y) = \min\left( 1, r(x,y)\right)$, where 
\[
r(x,y) := \frac{\pi(y)q(y,x)}{\pi(x)q(x,y)}
\]
is called the Hastings ratio (if $\pi(x)q(x,y) = 0$ then $\alpha(x,y) := 0$).  More generally any choice satisfying $\alpha(y,x) = r(x,y)\alpha(x,y)$ and $\alpha(x,y) \in [0,1]$ for all $x,y$ results in a Metropolis--Hastings transition kernel
\[
P(x,dy) = \alpha(x,y)Q(x,dy) + \left( 1 - \int \alpha(x,y)Q(x,dy) \right) \delta_x(dy)
\]
that is $\pi$-reversible.  The choice $\alpha(x,y) := \min(1, r(x,y))$ is almost always preferred as it minimises the asymptotic variance of ergodic averages \cite{peskun1973optimum,tierney1998note}.  We restrict attention to this choice throughout (though see Remark \ref{rem:general_alpha}).

A non-reversible algorithm that is still $\pi$-invariant can straightforwardly be constructed by cycling through different Metropolis--Hastings kernels in a manner that is not palindromic.  A simple example is to consider two candidate kernels $Q_1$ and $Q_2$ and respective Metropolis--Hastings kernels $P_1$ and $P_2$.  The kernel $P := P_1P_2$ is no longer reversible since $P^* = P_2P_1 \neq P$, but it is clearly still $\pi$-invariant.  One popular approach is to augment the state space to $\xi := (x,p)$, and define $P_1$ as a Metropolis--Hastings kernel with $Q_1(\xi, d\xi') := \delta_{T(\xi)}(d\xi')$ with $T$ an involution (meaning $T^{-1} = T$).  The kernel $P_2$ then only updates the \textit{momentum} $p$, leaving the \textit{position} $x$ unchanged.  Extensions of Peskun's theorem to non-reversible chains are given in \citep{maire2014comparison,andrieu2021peskun}.

Throughout we assume that $\pi$ is defined on $\mathbb{R}^d$.  Recall that a function $f: \mathbb{R}^d \to (0,\infty)$ is called \textit{lower semicontinuous} if for each $c \in \mathbb{R}$, the set $\{x \in \mathbb{R}^d : f(x) \le c\}$ is closed, or equivalently, $x_n\to x$ implies $\liminf_{n\to\infty} f(x_n) \ge f(x)$ (e.g. \cite[Section 2.1]{aliprantis2006infinite}). A function $g:\mathbb{R}^d \to (0,\infty)$ is called \textit{upper semicontinuous} if and only if $1/g(x)$ is lower semicontinuous.  A Markov transition kernel $P$ is called a $T$-chain if there is a distribution $\nu$ on $\mathbb{N}$ and a subtransition kernel $T$ such that $\int P^n(x,A)\nu(dn) \geq T(x,A)$, $T(x,\mathbb{R}^d) > 0$ for all $x$ and the function $f_A(x) := T(x,A)$ is lower semicontinuous for all $A \in \mathcal{B}(\mathbb{R}^d)$. In many cases $\nu = \delta_1$ is a reasonable choice.  The concept of a $T$-chain is introduced in Chapter 6 of \citep{meyn2012markov}, and allows straightforward characterisation of many topological and stability properties of $P$.  

A Markov chain with transition $P$ will be called geometrically ergodic if there is a $\pi$-a.e. finite $V \geq 1$, an $R < \infty$ and $\rho < 1$ such that for all $n \geq 1$ and any $x$ for which $V(x) < \infty$
\begin{equation}
\label{eq:geometric_ergodicity}
\|P^n(x,\cdot) - \pi\|_{V} \leq RV(x)\rho^n,
\end{equation}
where for a signed measure $\mu$ the $V$-norm $\|\mu\|_V := \sup_{|g|\leq V}\left|\int g(x)\mu(dx)\right|$.  It is well-known that among $\pi$-reversible chains geometric ergodicity is equivalent to the existence of a positive $L^2(\pi)$ spectral gap \citep{kontoyiannis2012geometric}.  The theory allows for $V$ to be infinite on a $\pi$-null set, however if an $x$ is chosen such that $V(x) = \infty$ then the bound ceases to be useful.  We therefore use \textit{geometric ergodicity} to refer exclusively to the case in which $V$ is finite for all $x$ in the support of $\pi$; this allows the useful equivalence in Proposition \ref{prop: MH is a T-chain}(ii), following Proposition 3.1 of \citep{roberts1996geometric}.

There are two classical settings in which the Metropolis--Hastings algorithm is well-known to perform poorly. The first is when $\alpha(x,y) \approx 0$ in some region of the state space, meaning almost all proposals are rejected.  It is shown in \cite{roberts1996geometric} that geometric ergodicity cannot then hold.  The other is when $\alpha(x,y) \approx 1$ for almost all proposals $y$ and $x$ sufficiently large, and $Q$ itself is not geometrically ergodic but meanders aimlessly.  This is the canonical `random walk behaviour' scenario, yet to the best of our knowledge no result directly characterising its relation to geometric ergodicity exists (though related results appear in \cite{jarner2003necessary}).  We provide one as Theorem \ref{them:non-geometric-ergodicity}, under mild regularity conditions on $\pi$ and $Q$ stated below.

The following assumption facilitates Proposition \ref{prop: MH is a T-chain} below, a generalisation of Theorem 2.2 in \cite{roberts1996geometric}, from which precise statements about $P$ will follow.

% \begin{lemma}
%     Let $f: \mathcal{X}\to [-\infty,+\infty]$ be a function on a topological space $\mathcal{X}$. Then:
%     \begin{itemize}
%         \item $f$ is lower semicontinuous if and only if for each $c \in \mathbb{R}$, the set $\{x \in \mathcal{X} : f(x) \le c\}$ is closed, or equivalently, $x_n\to x$ implies $\liminf_{n\to\infty} f(x_n) \ge f(x)$.
%         \item $f$ is upper semicontinuous if and only if for each $c \in \mathbb{R}$, the set $\{x \in \mathcal{X} : f(x) \ge c\}$ is closed, or equivalently, $x_n\to x$ implies $\limsup_{n\to\infty} f(x_n) \le f(x)$.
%     \end{itemize}
%     \label{lem:denfinition of semicontinuous}
% \end{lemma}

\begin{assumption}
    The distribution $\pi$ is defined on $\mathbb{R}^d$ and the following hold
    \begin{itemize}
        \item[(i)] $\pi$ has a density $\pi(x)$ that is upper semicontinuous and bounded away from 0 and $\infty$ on compact sets.
        \item[(ii)] $Q(x,\cdot)$ has a density $q(x,y)$, both $q(x,\cdot)$ and $q(\cdot,y)$ are lower semicontinuous for all $x$ and $y$, and there exists $\delta_q > 0$ such that for every $x \in \mathbb{R}^d$ \[
    \|x-y\| \leq \delta_q \implies q(x,y) > 0.
    \]
    \end{itemize}
    \label{ass:pi_q_regularities}
\end{assumption}

\begin{proposition}
\label{prop: MH is a T-chain}
    Let $\{X_n\}_{n\geq 0}$ be a Metropolis--Hastings Markov chain on $\mathbb{R}^d$ such that $\pi$ and $Q$ satisfy Assumption \ref{ass:pi_q_regularities}.  Then
    \begin{itemize}
        \item [(i)] $\{X_n\}_{n\geq 0}$ is a $\pi$-irreducible and aperiodic $T$-chain, and therefore all compact sets are small.
        \item [(ii)]  It is geometrically ergodic if and only if there exists a function $V: \mathbb{R}^d\to[1,\infty)$ such that 
        $$
        \limsup_{\|x\|\to\infty} \frac{PV(x)}{V(x)}<1.
        $$
    \end{itemize}
\end{proposition}

\begin{proof}
See Appendix \ref{app:proofs_MH_high_accept}.
\end{proof}

Our main result of this section is stated as Theorem \ref{them:non-geometric-ergodicity} below.  Intuitively, one might expect that if the acceptance rate of a Metropolis--Hastings algorithm approaches 1 in the tail region, then the transition kernel $P$ will behave similarly to $Q$ in this region. The tail behaviour of $Q$ should then determine the convergence rate of the algorithm. Consequently, if $Q$ fails to be geometrically ergodic and the acceptance rate approaches 1 in the tail, then $P$ should also fail to be geometrically ergodic.  We show below that this reasoning is essentially correct, but that a slightly stronger condition is required (condition (ii) below) than might be expected on the rate at which the acceptance rate tends to unity as $\|x\|\to \infty$.

\begin{theorem}
\label{them:non-geometric-ergodicity}
Let $\pi$ and $Q$ satisfy Assumption \ref{ass:pi_q_regularities}, and suppose in addition
\begin{enumerate}
    \item[(i)] $Q$ is not geometrically ergodic.
    \item[(ii)] For any $V : \mathbb{R}^d \to [1,\infty)$ such that $\int V(x)\pi(dx) < \infty$ it holds that 
    \[
    \lim_{\|x\|\to\infty} \int \frac{V(y)}{V(x)}(\alpha(x,y) - 1) Q(x,dy) = 0.
    \]
\end{enumerate}
Then the Metropolis--Hastings kernel $P$ is not geometrically ergodic.
\end{theorem}

\begin{proof}
First note that for any suitable $V$
\begin{align}
\label{eq:pv_expression}
    \frac{PV(x)}{V(x)}
    &= \int \frac{V(y)}{V(x)} \alpha(x,y)Q(x,dy) + \int \frac{V(y)}{V(x)} \delta_x(dy)\left(1- \int\alpha(x,y) Q(x,dy)\right)\nonumber\\
    &= \int \frac{V(y)}{V(x)} Q(x,dy) +\left(1- \int\alpha(x,y) Q(x,dy)\right)+D(x)
\end{align}
where $D(x) := \int V(y)(\alpha(x,y)-1)Q(x,dy)/V(x)$.  Under Assumption \ref{ass:pi_q_regularities} the extension of Proposition \ref{prop: MH is a T-chain}(ii) to the Markov chain produced by iterating $Q$ is immediate.   Condition (i) of the theorem statement therefore implies
\[
\limsup_{\|x\|\to\infty}\int \frac{V(y)}{V(x)}Q(x,dy) \geq 1.
\]
The second term in \eqref{eq:pv_expression} is non-negative (although by assumption it asymptotes to 0), and by condition (ii) of the theorem statement $\lim_{\|x\|\to\infty}D(x) = 0$. Combining gives
\begin{equation}
    \limsup_{\|x\|\to\infty}\frac{PV(x)}{V(x)} \geq \limsup_{\|x\|\to\infty}\int\frac{V(y)}{V(x)}Q(x,dy) \geq 1
\end{equation}
From this it is clear that $P$ cannot be geometrically ergodic by applying Proposition \ref{prop: MH is a T-chain}(ii).
\end{proof}

\begin{remark}
    When $Q$ satisfies a random walk type condition, similar results follow from \cite{jarner2003necessary}. Theorem \ref{them:non-geometric-ergodicity} also covers cases where this fails, such as the schemes in \cite{livingstone2021geometric,girolami2011riemann,roberts2002langevin,xifara2014langevin}.
\end{remark}

\begin{remark} \label{rem:general_alpha}
    The above can be straightforwardly extended to the setting in which $\alpha(x,y)$ is not chosen to be the $\min(1, r(x,y))$, provided that certain regularity conditions are still satisfied.  It may then be that $\lim_{\|x\|\to\infty}\int \alpha(x,y)Q(x,dy)  = c$ for some $c < 1$, which results in $P$ being a $(1-c)$-lazy version of $Q$ in the tails.
\end{remark}

It is tempting to conjecture that if $Q$ is not geometrically ergodic and $\lim_{\|x\|\to\infty} \int \alpha(x,y)\,Q(x,dy) = 1$ then $P$ will not be geometrically ergodic. This would provide a weaker sufficient criterion than Theorem \ref{them:non-geometric-ergodicity}(ii), since $\lim_{\|x\|\to\infty} \int \alpha(x,y)\,Q(x,dy) = 1$ is a special case of $\lim_{\|x\|\to\infty} \int V(y)/V(x)[\alpha(x,y)-1]\,Q(x,dy) = 0$ with $V(x) \equiv 1$. This conjecture is false, however, as demonstrated by the following counterexample.

\begin{proposition}
\label{prop: counterexample}
    Take $\pi(x)\propto \exp(-x^2)$ and set
    $$
    Q(x,dy) = [1- \varepsilon(x)] Q_1(x,dy)+\varepsilon(x) Q_2(x,dy),
    $$
    where $\varepsilon(x) = 1/(2+|x|)$, $Q_1(x,\cdot)$ denotes a $N(x/2,1)$ distribution, and $Q_2(x,\cdot)$ denotes a uniform distribution on $(\exp(x^2),\exp(x^2)+1)$.  Then $\lim_{|x|\to\infty} \int \alpha(x,y) Q(x,dy) = 1$ and $Q$ is not geometrically ergodic, but the Metropolis--Hastings kernel $P$ \emph{is} geometrically ergodic.
\end{proposition}

\begin{proof}
See Appendix \ref{app:proofs_MH_high_accept}.
\end{proof}

The intuitive explanation for Proposition \ref{prop: counterexample} is as follows. The candidate kernel $Q(x,dy)$ is a mixture of a well-behaved distribution $Q_1(x,dy)$ with probability $1-\varepsilon(x)$ and a large jump that occurs with probability $\varepsilon(x)$. In the tail region, nearly all proposals from $Q_1(x,dy)$ are accepted, while all large jumps are rejected. Since $\lim_{|x|\to\infty} \varepsilon(x) = 0$, almost all proposals in the tail come from $Q_1(x,dy)$, causing the average acceptance rate to approach 1. Although the large jump has low probability, when it occurs $V(y)$ becomes extremely large (where $y$ is the proposed state). This prevents $Q$ from being geometrically ergodic. The Metropolis--Hastings algorithm, however, rejects nearly all such large jumps, making the transition kernel $P$ behave as $Q_1$, which is geometrically ergodic.

A natural extension of Theorem \ref{them:non-geometric-ergodicity} is to consider the case in which $\pi$ has a density that becomes asymptotically flat in only some directions.  This is common in statistical models in the presence of near non-identifiability. As a simple example let $\pi(x) = \pi_1(x_1)\pi_2(x_2)$ and $Q(x,dy) := Q_1(x_1,dy_1)Q_2(x_2,dy_2)$ with a symmetric transition density.  If $\pi_1$ has polynomial tails then taking the sequence $x^{(n)} := (n, x_2)$ for $x_2$ fixed it holds that 
\[
\lim_{n\to\infty}\int \alpha(x^{(n)}, y)Q(x^{(n)}, dy) = \int \min\left(1, \frac{\pi_2(y_2)}{\pi_2(x_2)}\right)Q_2(x_2, dy_2).
\]
We do not therefore have that $\alpha \to 1$, but rather that it becomes asymptotically independent of $x_1$ in some directions. The $x_1$-coordinate process will then behave as an $x_2$-modulated lazy version of $Q_1$ in some areas of $\mathbb{R}^2$.  When $Q$ is a random walk kernel, for example, it can be shown that such chains will not be geometrically ergodic using the results of \cite{jarner2003necessary}.  

The previous paragraph connects with the popular idea of acceleration in sampling and optimisation. In the sampling context, take $\pi$ to have an $m$-strongly log-concave density $\pi(x)$ for which $\nabla \log\pi(x)$ is gradient-Lipschitz with constant $M$. A condition number $\kappa := M/m \geq 1$ can then be associated with $\pi$, and often connected with the mixing time of a Markov chain algorithm targeting $\pi$. Often the mixing time is linear in $\kappa$, but for some \textit{accelerated} algorithms it is reduced to $\sqrt{\kappa}$.  Large $\kappa$ typically means $\pi$ is flatter in some directions than others, and mixing is determined by how well these flatter directions are explored. The square root improvement follows directly from the fact that accelerated chains move ballistically rather than diffusively in these directions (see e.g. \cite{oberdorster2025accelerated}).  Linear pre-conditioning can reduce the condition number and thereby increase sampling efficiency \cite{hird2025quantifying}. Similarly, applying a nonlinear transformation to induce faster-decaying tails can recover geometric ergodicity \cite{johnson2012variable}.

%% OLD Stuff %%

% \begin{lemma}
%     Let $f: X\to [-\infty,+\infty]$ be a function on a topological space. Then:
%     \begin{itemize}
%         \item [] f is lower semicontinuous if and only if for each $c \in \mathbb{R}$ the set $\{x \in \mathcal{X} : f(x) \le c\}$ is closed or $x_n\to x\Rightarrow \lim_{n\to\infty}\inf f(x_n) \ge f(x)$
%         \item[] f is upper semicontinuous if and only if for each $c \in  \mathbb{R}$ the set $\{x \in \mathcal{X} : f(x) \ge c\}$ is closed or $x_n\to x\Rightarrow \lim_{n\to\infty}\sup f(x_n) \le f(x)$
%     \end{itemize}
%     \label{lem:denfinition of semicontinuous}
% \end{lemma}

% \begin{assumption}
%     Assume the following:
%     \begin{itemize}
%         \item[(i)] $\pi(dx) := \pi(x)dx$ with the density $\pi(x)$ is upper semicontinuous.
%         \item[(ii)] $Q(x,dy):= q(x,y) dy$ with the density $q(x,y)$ is  lower semicontinuous.
%     \end{itemize}
%     \label{ass:pi_q_regularities}
% \end{assumption}


% \begin{proposition}
%     Let $\{X_n\}_{n\geq 0}$ be a Metropolis--Hastings chain on $\mathbb{R}^d$ with a target distribution $\pi$, proposal kernel $Q(x,dy) := q(x, y)dy$, and acceptance probability $\alpha(x, y)$, the transition kernel $P(x,dy)$, if the target distribution and proposal kernel satisfies the assumption \ref{ass:pi_q_regularities}, then we have:
%     \begin{itemize}
%         \item [(i)] The chain $\{X_n\}_{n\geq 0}$ is a $\pi$-irreducible and aperiodic $T$-chains
%         \item [(ii)] Every every compact set $C \subset \mathbb{R}^d$ is small. The sufficient and necessary condition for the geometrical ergodicity of the chain is, there exists one function $V: \mathbb{R}^d\rightarrow[1,+\infty)$ such that 
%         $$
%         \lim_{|x|\to\infty}\sup \frac{PV(x)}{V(x)}<1
%         $$
%     \end{itemize}
% \end{proposition}

% \begin{proof}
%     Firstly we proof the chain is a $T$-chain. According to the Chapter 6 in \cite{meyn2012markov}, we know the chain is a $T$-chain if the transition kernel satisfies
%     $$
%     P(x,A) \ge T(x,A),\quad x\in \mathcal{X}, A\in \mathcal{B}(\mathcal{X})
%     $$
%     with $T(\cdot,A)$ is a semicontinuous function for any $A\in \mathcal{B}(\mathcal{X})$. 
%     Then we know:
%     $$
%     P(x,A) = \int_A \alpha(x,y)q(x,y)dy+r(x)\delta_x(dy) \ge\int_A \alpha(x,y)q(x,y)dy:=T(x,A)
%     $$
% Then we want to prove $T(x,A)$ is lower semicontinuous for any $A \in \mathcal{B}(X)$. For achieving this, we firstly denote 
% \begin{align}
%     T_y(x) &:= \alpha(x,y)q(x,y)
%     = \min\{q(x,y),\frac{\pi(y)q(y,x)}{\pi(x)}\}
% \end{align}
%     because $T_y(x)$ is a function with respect to $x$, so here we just treat $y$ as a constant. We know $\pi(x)$ is upper semicontinuous, based on the definition of upper/ lower semicontinuous, we can easily get that $1/\pi(x)$ is lower semicontinuous. Then based on Lemma \ref{lem:denfinition of semicontinuous}, we know, for any sequence $\{x_n\}\to x$, we have 
%     $$
%     \lim_{n\to\infty}\inf q(y,x_n)\ge q(y,x)\ge 0; 
%     $$
%     $$
%     \lim_{n\to\infty}\inf 1/\pi(x_n)\ge 1/\pi(x)\ge 0
%     $$
%     because $q(y,x_n),1/\pi(x_n)$ is not negative so we can get,
%     $$
%     \lim_{n\to\infty}\inf \frac{q(y,x_n)}{\pi(x_n)}\ge ( \lim_{n\to\infty}\inf q(y,x_n)) \cdot ( \lim_{n\to\infty}\inf 1/\pi(x_n))\ge \frac{q(y,x)}{\pi(x)}
%     $$
%     and $y$ is a constant so $\pi(y) q(y,x)/\pi(x)$ is a lower semicontinuous function with respect to $x$. And we can prove that if $g,h$ are both lower semicontinuous functions, then $\min\{g,h\}$ is also semicontinuous function. Because for each $c \in \mathbb{R}$ $\{x\in X:\min\{g(x),h(x)\}\le c\} =\{x\in X:g(x)\le c\} \cup \{x\in X:h(x)\le c\} $, because $h(x)$ and $g(x)$ are both semicontinuous, then the two sets are closed, so the union of two closed sets are also closed. So $\min\{g,h\}$ is lower semicontinuous.

%     Then here we get $q(x,y)$ and $\pi(y)q(y,x)/\pi(x)$ are both lower semicontinuous function, so $T_y(x)$ is also semicontinuous function. Then $T(x,A) = \int_A T_y(x) dy$, then for any sequence $x_n\to x$, we have
%    \begin{align}
%          \lim_{n\to\infty}\inf T(x_n,A) &= \lim_{n\to\infty}\inf \int_A T_y(x_n)dy = \lim_{n\to\infty}\inf  \int T_y(x_n)\mathbb{I}_A(y)dy
%    \end{align}
%   because $T_y(x_n)\ge 0$, based on the Fatou's Lemma, we know
%   $$
%   \lim_{n\to\infty}\inf  \int T_y(x_n)\mathbb{I}_A(y)dy\ge   \int \lim_{n\to\infty}\inf T_y(x_n)\mathbb{I}_A(y)dy
%   $$
%   And based on the $T_y(x)$ is a lower semicontinuous function, we know $\lim_{n\to\infty}\inf T_y(x_n)\ge T_y(x)$, Thus
%   \begin{equation}
%       \lim_{n\to\infty}\inf T(x_n,A)\ge \int \lim_{n\to\infty}\inf T_y(x_n)\mathbb{I}_A(y)dy\ge \int_A T_y(A)  = T(x,A)
%   \end{equation}
%  So, we know $T(x,A)$ is a lower semicontinuous function for any $A\subset \mathcal{B}(X)$. So the chain is a $T$-chain.

%  Then we prove this chain is $\pi$-irreducible and aperiodic.

%  Then base on the theorem 5.5.7 and theorem 6.2.5 in \cite{meyn2012markov}, we know if the chain is $\pi$-irreducible and aperiodic, then every petite is small and if a chain is a $\pi-$ irreducible $T$-chain then every compact set is petite. 

%  So we can conclude if a Markov chain is a $\pi-$irreducible and aperiodic $T$-chain, then the compact set is small. And based on the Theorem 1.4 in \cite{mengersen1996rates}, we know if a chain is geometrically ergodic, it is equivalent that there is a function $V\ge 1$, finite at least for one $x$ and a small set $C$ such that fo some $\lambda<1$, drift condition
%  $b<\infty$, 
%  $$
%  PV\le \lambda V + b\mathbb{I}_C
%  $$
%  is satisfied.
% The condition $\limsup_{|x| \to \infty} PV(x)/V(x) < 1$ implies that there exists some $\lambda < 1$ and a radius $M$ such that for all $|x| > M$, $PV(x)/V(x) \leq \lambda$.Define the compact set $C = \{x : |x| \leq M\}$.We know the compact set is also a small set. So this condition is equivalent to the drift condition.
    
% \end{proof}

% \newpage


% \begin{proposition}
% \label{prop: MH is a T-chain}
% Let $\{X_n\}_{n\geq 0}$ be a Metropolis--Hastings chain on $\mathbb{R}^d$ with a target distribution $\pi$, proposal kernel $Q(x,dy) := q(x, y)dy$, and acceptance probability $\alpha(x, y)$.  Assume the following:
% \begin{enumerate}
%     \item $\pi$ has a density $\pi(x)$ that is bounded away from $0$ and $\infty$ on any compact set $C \subset \mathbb{R}^d$.
%     \item For any compact set $C \subset \mathbb{R}^d$, there exist $\delta_C > 0$ and $\epsilon_C > 0$ such that:
%     \[
%     q(x, y) \geq \epsilon_C \quad \text{whenever } x \in C \text{ and } \|x - y\| \leq \delta_C.
%     \]
%     \item For any compact set $C$, there exists $\eta_C > 0$ such that:
%     \[
%     \alpha(x, y) \geq \eta_C \quad \text{whenever } x \in C \text{ and } \|x - y\| \leq \delta_C
%     \]
% \end{enumerate}
% Then the chain is $\pi$-irreducible, aperiodic, and every compact set $C \subset \mathbb{R}^d$ is small.
% \end{proposition}

% \begin{proof}
% Define the sub-stochastic density $p(x, y) := q(x, y)\alpha(x, y)$.  A set $C$ is small if there exists a measure $\nu$, an integer $m \ge 1$, and $\beta > 0$ such that $P^m(x, A) \ge \beta \nu(A)$ for all $x \in C$. Consider any compact $C$. By conditions (2) and (3):
% \[
% p(x, y) = q(x, y)\alpha(x, y) \geq \epsilon_C \eta_C > 0
% \]
% whenever $x \in C$ and $|x - y| \leq \delta_C$.
% %The function $p(x, y)$ is strictly bounded away from zero in a $\delta_C$-neighborhood of the diagonal on $C \times \mathbb{R}^d$. This implies that $p(\cdot, \cdot)$ is lower semicontinuous on $C \times \mathbb{R}^d$, meaning that $P$ is a $T$-chain, and therefore provided that it can be shown to be $\pi$-irreducible then every compact set is small.
% Fix $x_0 \in C$. Let $B_{\delta_C/2}(x_0) \subset C$ be a Euclidean ball centered at $x_0$ with radius $\delta_C/2$. For any $x \in B_{\delta_C/2}(x_0)$ and $y \in B_{\delta_C/2}(x_0)$, $|x - y| \leq \delta_C$. Thus, for all $x \in B_{\delta_C/2}(x_0)$:
% \[
% P(x, A) \ge \int_{A \cap B_{\delta_C/2}(x_0)} \epsilon_C \eta_C \, dy = \beta \nu(A),
% \]
% where $\nu$ is the Lebesgue measure restricted to $B_{\delta_C/2}(x_0)$. Thus, sufficiently small compact sets are 1-small. %Under the local positivity and continuity conditions established above, every point $x \in \mathbb{R}^d$ possesses a neighborhood that is $1$-small. 
% By the Heine--Borel theorem, any compact set is covered by a finite sub-cover of these neighborhoods, which implies that any compact set is petite. Given that the chain is a $\pi$-irreducible and aperiodic, by Theorem 5.5.7 of \cite{meyn2009markov} implies that every compact set is also small.

% To establish $\pi$-irreducibility, note that for any $(x,A) \in \mathbb{R}^d \times \mathcal{B}(\mathbb{R}^d)$ we can choose an arbitrary compact $A^* \subset A$ and then define another compact and connected set $C(x,A^*)$ containing both $x$ and $A^*$.  We can then connect $x$ and any point in $A^*$ with a finite sequence of overlapping Euclidean balls, each of radius $\delta_{C(x,y)}/2$.  This implies that there exists a finite $m$ such that $P^m(x,A)>0$. The result follows since $x$ and $A$ are arbitrary.

% %Since $\pi$ is everywhere positive, $\pi(A)>0$ implies $A$ has positive Lebesgue measure. 

% %Condition (2) implies that from any $x \in C$, the chain transitions with positive density to any open ball within distance $\delta_C$. Since $\mathbb{R}^d$ is connected, any two points can be connected by a finite sequence of overlapping $\delta$-balls. Thus, there exists a path of densities $p(x_0, x_1) \dots p(x_{n-1}, z) > 0$, ensuring $\pi$-irreducibility.

% From Condition (2), $p(x, y) > 0$ for $|x-y| \le \delta_C$. This allows arbitrarily small moves and overlap between $P(x, \cdot)$ and $P(z, \cdot)$ for $x, z$ close to each other, which prevents rigid cyclic behavior and ensures aperiodicity.

% {\color{blue} Leave this part for Sam}
% \end{proof}

% \begin{assumption} \label{ass:pi_q_regularity}
% Assume the following:
% \begin{enumerate}
%     \item[(i)] $\pi(dx) := \pi(x)dx$ with the density $\pi(x)$ bounded away from $0$ and $\infty$ on compact sets.
%     \item[(ii)] $Q(x,dy) := q(x,y)dy$ and that there exists $\delta_q > 0$ and $\epsilon_q > 0$ such that for every $x \in \mathbb{R}^d$
%     \[
%     \|x-y\| \leq \delta_q \implies q(x,y) \geq \epsilon_q
%     \]
% \end{enumerate}
% \end{assumption}
% {\color{red} 
% \begin{remark}
% The assumptions fix $\epsilon_q$ and $\delta_q$ for all $x$ meaning that there is some restriction imposed by them.  I don't think this is actually needed, I think we might just need $q(x,y) > 0$ or we can weaken the condition in some other way, but we need to check this...
% \end{remark}
% }
% {\color{blue}
% If this assumption is only used to state the equivalent condition for geometric ergodicity, then based on Theorem 1.4 of \cite{mengersen1996rates}, we may only need to require that the chain $\Phi$ is $\pi$-irreducible and aperiodic. This means we may only need to impose the conditions that $q(x,y)>0$ for all $x,y$ and that $q(x,y)$ is a continuous function.
% }
\section{Random \& Guided Walk Metropolis on $\mathbb{R}$}
\label{sec: RWM and GWM on 1D R}

We now turn to the random walk and guided walk Metropolis on $\mathbb{R}$. The (Gaussian) random walk Metropolis on $\mathbb{R}$ is the case of Metropolis--Hastings in which $Q(x,\cdot)$ is $N(x, \epsilon^2)$. The candidate density $q(x,y)$ is then symmetric in its arguments, so the Hastings ratio reduces to $r(x,y) = \pi(y)/\pi(x)$. The algorithm is well-known and well-studied (e.g. \cite{sherlock2010random}).

The guided walk Metropolis was introduced in \cite{GWM_intro} as a simple way to incorporate momentum into a sampling algorithm.  The state space is augmented to be $\mathbb{R} \times \{-1,+1\}$ with the new state variable $\xi := (x,p)$ such that $p \in \{-1,+1\}$ enforces a direction to the transitions of the Markov chain.  The algorithm produces a Markov transition from the current state $(X_i,P_i)$ by generating the proposal
\[
Y = X_i + \epsilon P_i|Z|
\]
for some $Z \sim N(0,1)$, and then setting 
\[
(X_{i+1},P_{i+1}) 
= 
\begin{cases}
    (Y,P_i), & \text{with probability } \alpha(X_i,Y) \\
    (X_i,-P_i), & \text{otherwise},
\end{cases}
\]
where $\alpha(X_i,Y) = \min(1, \pi(Y)/\pi(X_i))$ as in the random walk Metropolis.  Iterating produces a Markov chain on $\mathbb{R} \times \{-1,+1\}$ with invariant distribution $\mu(d\xi) = \pi(dx)\nu(dp)$, where $\nu$ denotes the uniform distribution on $\{-1,+1\}$.  The guided walk transition can be viewed as the two-skeleton of a two-cycle Metropolis--Hastings kernel $P_1P_2$, where $P_1$ has candidate kernel $Q_1(\xi,d\xi') = \int \delta_{(x+\epsilon p|z|,-p)}(dx',dp')\phi(z)dz$, with $\phi$ denoting the standard Normal density on $\mathbb{R}$, and $Q_2$ consists of keeping $x$ fixed and setting $p \to -p$, a candidate transition that will always be accepted owing to the symmetry of $\nu$.  Further discussion on the general construction of such deterministic transitions is given in \cite{tierney1998note,andrieu2020general, glatt2023accept}.

The key advantage of the guided walk is that if a proposal $Y$ is accepted then the chain continues moving in the same direction.  If the chain is in a region for which $\pi(x)$ is relatively flat, therefore, this will promote ballistic motion when the random walk Metropolis will instead behave diffusively.  In \cite{andrieu2021peskun} it is shown that if $P$ denotes the guided walk transition then $(P+P^*)/2$ is the random walk Metropolis (or rather an augmented version on $\mathbb{R} \times \{-1,+1\}$ for which the momentum is completely refreshed at every iteration).  When a proposal is rejected, however, the direction of travel in the guided walk is changed, and so in the presence of a non-trivial $\alpha$ the degree of benefit over the random walk Metropolis is not so clear.  We study both the polynomial-tailed case, in which $\pi$ becomes flat for large $|x|$, and the strictly log-concave case, in which $-\log \pi(x)$ grows faster than linearly in $|x|$.

% \begin{align*}
%     y &= x_n + p_n|z|, \quad Z \sim N(0, \epsilon^2)
%     \alpha(x_n, y) &= \min\left\{1,\frac{\pi(y)}{\pi(x_n)}\right\}\\
%     (x_{n+1},p_{n+1}) &= \begin{cases}
%     (y,p_n), & \text{with probability } \alpha \\
%     (x_n,-p_n), & \text{with probability } 1-\alpha
%     \end{cases}
% \end{align*}



% In this section, we will investigate whether the classic random-walk Metropolis (RWM) and other Metropolis algorithms (such as the guided Walk Metropolis (GWM) \citep{GWM_intro}) which claim to achieve `diffusive-to-ballistic' speed-up by suppressing the random-walk behaviour behave in a similar manner at the the tail for the target distribution. Ideally these new introduced methods should achieve a faster convergence, however, we will show that the target distribution $\pi$ affects the mixing behaviour between these two MCMC algorithms. We claim that, when the target density is light-tailed, RWM and a lazy-$1/2$ GWM do not behave differently at the tail, while such ballistic speed-up in sampling is observed for heavy tailed target distributions.


% \subsection{Background: RWM \& GWM}
% \label{subsec: background: RWM, GWM}
% We first give a brief recap of what is meant by RWM and GWM, also its lazy-$1/2$ GWM (shorten as LGWM). Define the random walk Metropolis (RWM) with Gaussian noise $z_n \sim N(0,1)$ and fixed stepsize $\epsilon$:
% \begin{equation}
%     X_{n+1} = X_n + \epsilon z_{n+1} \implies q(X_{n+1} \mid X_n = x) = N(x; X_n, \epsilon^2)
% \end{equation}
% and write the 1-step RWM transition kernel as $P_{\text{RWM}}(x, \cdot)$ if we start from $X_n = x$. Based on RWM, this kernel is expressed as
% \begin{equation}
% \label{eqn: RWM transition kernel}
%     P_{\text{RWM}}(x, dy) = \alpha(x, y)q(x, dy) + (1 - \rho(x))\delta_x(dy)
% \end{equation}
% where $\rho(x) = \int_{\mathcal{X}} \alpha(x,y)q(x,dy)$, with the proposal $q(x,y) = N(y;x,\epsilon^2) = \frac{1}{\epsilon\sqrt{2\pi}}\exp\left(-\frac{1}{2\epsilon^2}(y-x)^2\right)$ and the acceptance rate $\alpha(x,y) = \min\left(1, \frac{\pi(y)}{\pi(x)}\right) = \min \left(1, \exp \left(\frac{1}{2}(x^2 - y^2) \right)\right)$.

% Based on RWM algorithm this kernel is expressed as
% \begin{equation}
% \label{eqn: RWM transition kernel}
%     P_{\text{RWM}}(x, dy) = \alpha(x, y)q(x, dy) + (1 - \rho(x))\delta_x(dy)
% \end{equation}
% where $\rho(x) = \int_{\mathcal{X}} \alpha(x,y)q(x,dy)$ and
% \begin{itemize}
%     \item The proposal density:
%     \begin{equation}
%     \label{eqn: proposal density q}
%         q(x,y) = N(y;x,\epsilon^2) = \frac{1}{\epsilon\sqrt{2\pi}}\exp\left(-\frac{1}{2\epsilon^2}(y-x)^2\right)
%     \end{equation}

%     \item The MH acceptance rate/probability (with this symmetric proposal $q$):
%     \begin{equation}
%         \alpha(x,y) = \min\left(1, \frac{\pi(y)}{\pi(x)}\right) = \min \left(1, \exp \left(\frac{1}{2}(x^2 - y^2) \right)\right)
%     \end{equation}
% \end{itemize}

% Meanwhile, \cite{GWM_intro} introduces the 1D guided walk metropolis (GWM), written as
% \begin{align*}
%     y &= x_n + \epsilon p_n|z|, \quad z\sim N(0, 1)\\
%     \alpha(x_n, y) &= \min\left\{1,\frac{\pi(y)}{\pi(x_n)}\right\}\\
%     (x_{n+1},p_{n+1}) &= \begin{cases}
%     (y,p_n), & \text{with probability } \alpha \\
%     (x_n,-p_n), & \text{with probability } 1-\alpha
%     \end{cases}
% \end{align*}
% To summarise, given the target density $\pi$ and the current/starting position-momentum state $(x,p)$, with the state space be $\mathcal{X}\times \{-1, 1\}$, we write our transition probability from $(x, p)$ to $A \times \{p'\}$, or in short as $A \times p'$, where $A\subseteq \mathcal{X}$ and $p' \in \{-1, 1\}$. GWM introduces two possible transition cases of the position-momentum pair that
% \begin{enumerate}
%     \item Transition position and keep momentum: $(x, p) \to (y, p)$ with probability $\alpha(x,p,y)$.
%     \item Keep position and transition momentum: $(x,p) \to (x, -p)$ otherwise.
% \end{enumerate}
% Indeed, we can write the proposal density for GWM of the transition position given the starting position-momentum state $(x,p)$ in terms of
% \begin{equation}
% \label{eqn: GWM position proposal}
%     q(y \mid x,p) = q(x,p,y) = \frac{2}{\epsilon \sqrt{2\pi}}\exp\left(-\frac{1}{2\epsilon}(y-x)^2\right)\mathbf{1}(p(y-x)\geq 0)
% \end{equation}
% Note that here we replace $\sigma$ with $\epsilon$ as the standard deviation of the Gaussian noise. \eqref{eqn: GWM position proposal} matches the algorithm transition setting $y=x+\epsilon |z|, z\sim N(0,1)$ that we will apply the right branch of $\text{Half-Normal}(x, \epsilon^2)$ if $p=1$ (so $y>x$ resulting $p(y-x)\geq 0$), and vice versa. i.e. for $p\in \{-1,1\}$:
% \begin{align}
% \label{eqn: GWM proposal at p=-1, half normal left branch}
%     q(x,-1,y) &= \frac{2}{\epsilon \sqrt{2\pi}}\exp\left(-\frac{1}{2\epsilon}(y-x)^2\right)\mathbf{1}(y-x < 0) = p_1(x, y)\\
% \label{eqn: GWM proposal at p=1, half normal right branch}
%     q(x,1,y) &= \frac{2}{\epsilon \sqrt{2\pi}}\exp\left(-\frac{1}{2\epsilon}(y-x)^2\right)\mathbf{1}(y-x\geq 0) = p_2(x, y)
% \end{align}
% and $p_1, p_2$ is the density for the left/right branch of $\text{Half-Normal}(x, \epsilon^2)$ respectively. Note that the proposal \eqref{eqn: GWM position proposal} is symmetric with $q(x,p,y) = q(y,-p,x)$, resulting the acceptance probability as
% \begin{equation}
% \label{eqn: GWM acceptance rate, general form}
%     \alpha(x, p, y) = \min\left(1, \frac{\pi(y)}{\pi(x)}\right) = \min \left(1, \frac{\pi(x + p\epsilon |z|)}{\pi(x)} \right)
% \end{equation}
% With all above settings, we write the law of first-step transition for GWM from $x$ as
% \begin{equation*}
% \begin{aligned}
%     &P_{\text{GWM}}((x,p), A\times p') = P(\text{case 1 and $y\in A, p'=p$}) + P(\text{case 2 and $x\in A, p'=-p$})\\
%     &= \mathbf{1}(p'=p)\int_A \alpha(x,p,y) q(x,p,dy) + \mathbf{1}(p'=-p)\mathbf{1}(x\in A) \int_{\mathcal{X}}(1 - \alpha(x,p,y))q(x,p,dy)
% \end{aligned}
% \end{equation*}
% which is equivalently to write as the transition kernel for GWM
% \begin{equation}
% \label{eqn: GWM transition kernel}
%     P_{\text{GWM}}((x,p), (dy, p')) = \alpha(x,p,y)q(x,p,dy)\mathbf{1}(p'=p) + (1-\rho(x,p))\mathbf{1}(p'=-p)\delta_x(dy)
% \end{equation}
% where
% \begin{equation}
% \label{eqn: GWM transition kernel rho(x,p) general form}
%     \rho(x,p) = \int_\mathcal{X} \alpha(x,p,y)q(x,p,dy)
% \end{equation}
% We then can carry a marginalisation on $p'$, the newly transitioned momentum, to derive the GWM that is only interested on the state space related to positions. The marginalised form is written as
% \begin{equation}
% \label{eqn: marginalised GWM}
% \begin{aligned}
%     P((x,p), dy) &= \sum_{p'=\pm p}P_{\text{GWM}}((x,p), (dy, p'))\\
%     &= \alpha(x,p,y)q(x,p,dy) + (1-\rho(x,p))\delta_x(dy)
% \end{aligned}
% \end{equation}
% Due to the presence of auxiliary momentum variable $p$, it does not enable us to directly evaluate the TV distance between GWM and 1D Metropolis-Hastings based transition kernels (only related to starting position $x$, e.g. RWM). This invokes us to assign an initial $p=p(x) \in \{-1,1\}$ relating to $x$, so as to result a law of transition that is only related to the starting position and drifted position. Intuitively, a reasonable choice for such initial $p$ can be $p = -\operatorname{sgn}(x)$ to match the GWM behaviour, that under a symmetric target density (centered at 0), we will have high tendency to move towards 0, meaning if we are currently at the left ($x<0$)/right ($x\geq 0$) w.r.t the mode 0, we have higher tendency to have a right/left drift towards the high density region respectively, showing an opposite direction of moving w.r.t the current position $x$, which is precisely $-\operatorname{sgn}(x)$. The law of 1-step transition for GWM with $p=-\operatorname{sgn}(x)$ as the initial setting is then
% \begin{equation}
% \label{eqn: GWM law of transition, 1-step, initial p=-sgn(x)}
%     \mathcal{L}_{\text{GWM}}(x, dy) = \alpha(x,-\operatorname{sgn}(x),y)q(x,-\operatorname{sgn}(x),dy) + (1-\rho(x,-\operatorname{sgn}(x)))\delta_x(dy)
% \end{equation}

% \begin{definition}[Lazy-$1/2$ GWM (LGWM)]
% \label{def: LGWM}
%     We define the lazy version of GWM with a `lazy ratio' $1/2$ with law of 1-step transition as
%     \begin{equation}
%         \mathcal{L}_{\text{LGWM}}(x, dy) = \frac{1}{2}\mathcal{L}_{\text{GWM}}(x,dy)+\frac{1}{2}\delta_x(dy)
% \end{equation}
% \end{definition}

\subsection{Polynomial tails}
\label{subsec:heavy_tailed_target}

In this section we assume $\pi$ satisfies the following condition, which is representative of many heavy-tailed distributions on $\mathbb{R}$.
\begin{assumption}
\label{ass:heavy-tailed}
Assumption \ref{ass:pi_q_regularities}(i) holds and in addition there exist constants $r>0$, $K>0$, and $C_0>0$ such that for $|x|>K$,
$$
\pi(x) = \frac{C_0}{|x|^{1+r}}.
$$  
\end{assumption}

A Markov chain is called polynomially ergodic if there is a constant $\beta > 0$ and real-valued functions $V_\beta \geq 1$, $M > 0$  such that for all $n \geq 1$
\begin{equation}
\label{eq:polynomial_ergodicity}
\|P^n(x,\cdot) - \pi\|_{V_\beta} \leq M(x) n^{-\beta}.
\end{equation}
Note that in this case $V_\beta$ depends on $\beta$, and in fact there is an inherent trade-off between the rate of convergence and its choice \citep{jarner2007convergence}. We therefore focus on the total variation norm, i.e. choosing $V_\beta \equiv 1$.  The largest choice of $\beta$ for which \eqref{eq:polynomial_ergodicity} holds with $V_\beta \equiv 1$ is called the \textit{polynomial rate}.

Under Assumption \ref{ass:heavy-tailed} the following result is proven in \citep{jarner2007convergence} about the random walk Metropolis.

\begin{proposition}[Jarner \& Roberts, 2007]
    Under Assumption \ref{ass:heavy-tailed} the (Gaussian) random walk Metropolis is polynomially ergodic with rate $r/2$.
\end{proposition}

For the guided walk Metropolis, by contrast, we establish a faster polynomial rate of convergence.

\begin{proposition}
    Under Assumption \ref{ass:heavy-tailed} the guided walk Metropolis is polynomially ergodic with rate $r$.
\end{proposition}

\begin{proof}
    In Lemma \ref{lem:lowerbound_polynomialtails} we show that the polynomial rate is at least $r$, and in Lemma \ref{lem:upperbound_polynomialtails} we show that it is $\leq r$.  Both results can be found in Appendix \ref{app:proofs_gwm_polynomialtails}.
\end{proof}

The proof follows a similar argument to that used in \citep{vasdekis2022note} to study a related continuous-time process.  The result shows that in this setting a doubling in rate is achieved.  The guided walk Metropolis is a \textit{lifted} version of the random walk Metropolis, and it is known that when the state space is finite or compact a lifted Markov chain must have mixing time at least the square root of the original chain, which translates to at most a doubling of the rate \cite{chen1999lifting,ramanan2018bounds}.


\subsection{Strictly log-concave tails}
\label{subsec:light_tailed_target}

We now consider the case of $\pi$ having a density and write $U(x) := -\log\pi(x)$ to denote the associated \textit{potential} function.  Assume the following strict convexity and super-linear growth conditions.
\begin{assumption}
\label{ass:strict_convexity}
    The target distribution $\pi$ has a Lebesgue density $\pi(x) \propto e^{-U(x)}$ on $\mathbb{R}$ with $U \in C^1(\mathbb{R})$, $U'$ is strictly monotonically increasing and
    \begin{equation*}
    \lim_{|x| \to \infty}\frac{U(x)}{|x|} = \infty.
    \end{equation*}
\end{assumption} 
It follows straightforwardly from Assumption \ref{ass:strict_convexity} that $U'(x) \to \infty$ as $x \to \infty$ and $U'(x) \to -\infty$ as $x \to -\infty$.
% We claim $U'(x) \to \infty$ when $x \to \infty$, and it can be shown by contradiction: with convexity, assume $U'(x) \in (0, M), \forall x \in \mathbb{R}$ is bounded for some positive, finite $M$. Then for $0<p<x$ we have
% \begin{equation*}
% \begin{aligned}
%     U(x) - U(p) = \int_p^x U'(t)dt \leq \int_p^x M dt = M(x-p) &\implies U(x) \leq Mx + U(p)\\
%     &\implies \lim_{x\to \infty}\frac{U(x)}{x} \leq M
% \end{aligned}
% \end{equation*}
% which contradicts with our setting of super-linear growth, hence by contradiction we show $U'(x) \to \infty$ when $x \to \infty$. By the similar thinking we can also show $U'(x)\to -\infty$ when $x \to -\infty$.

% \subsubsection{One-step law of transition analysis}
% \label{subsubsec: first-step RWM and LGWM}
% We first analyse whether the law of the one-step transitions between the RWM and the LGWM behave similarly in the tails in the case $\pi(x)$ is a general light-tailed distribution satisfying (i)-(iii). We claim the following proposition:

% \begin{proposition}[Equivalence of one-step transition]
% \label{prop:1step_coupling}
% Denote $\mathcal{L}_{\text{RWM}}(x, \cdot)$ and $\mathcal{L}_{\text{LGWM}}(x, \cdot)$ as the law of first-step transition of RWM and LGWM starting at $x$ respectively, when the target density is general light-tailed distribution, we have the convergence of TV distance in terms of
% \begin{equation}
% \label{eqn: TV convergence statement, LGWM and RWM}
%     \lim_{|x| \to \infty} \|\mathcal{L}_{\text{RWM}}(x, \cdot) - \mathcal{L}_{\text{LGWM}}(x, \cdot)\|_{TV} = 0
% \end{equation}
% \end{proposition}

% \begin{proof}

% We couple the RWM and the LGWM proposals synchronously in the following way. Given the starting position $X_0 = x$ for RWM and $(W_0, P_0) = (x,p)$ for LGWM where $p := -\operatorname{sgn}(x)$, draw $z \sim N(0,1)$ and set the RWM proposal to be $x' = x + \epsilon z$ and the LGWM proposal to be $(w',p') = (x + \epsilon p|z|, p)$. Note that $\mathbb{P}P(pz > 0) = 1/2$ with deterministic $p$, corresponding to the lazy ratio 1/2. If $pz > 0$ note that $x' = w'$.  In this instance set $X_1 = W_1 = x'$ with probability $\alpha(x,x')$, otherwise let $X_1 = x$ and $(W_1, P_1) = (x, -p)$. If instead $pz \leq 0$ then $x' \neq w'$, but upon rejection $X_1 = W_1 = x$.  The next states for the position variable will therefore be equal if either $pz > 0$ and the proposals are both accepted, or in the case that both proposals are rejected.  Since in the former case $\alpha(x,x') = \alpha(x,w')$ then using a common $\mathrm{Unif}[0,1]$ draw to make the acceptance decision ensures that when $pz > 0$ then $X_1 = W_1$. To ensure that the laws become equal in general, we must show that if $pz \leq 0$ then the acceptance probability $\to 0$ as $x \to \infty$, and such probability is easily $\int_{x}^{\infty}\alpha(x,y)q(x, dy)/2$ (or $\int_{-\infty}^{x}\alpha(x,y)q(x, dy)/2$ if we consider the tail $(-\infty, x]$, with $x\to -\infty$), that converges to 0 by Lemma \ref{lem:convergence_alpha_light_tails},  which establishes the result using the coupling inequality for TV distance.

% \begin{algorithm}[ht]
% \caption{Coupling construction algorithm, RWM and LGWM, for $x>0$}\label{alg: coupling construction, RWM and LGWM}
% \begin{algorithmic}[1]

% \State{Given the starting position $x$, starting momentum $p=-\operatorname{sgn}(x)$.}

% \State{Sample $z \sim N(0,1)$, compute $x' = x + \epsilon z$ and $w' = x + \epsilon p|z|$.}


% \If{$z \leq 0$}
% \Comment{for this case we have $x'=w'=x-\epsilon |z|$, with $x>0 \implies p=-1$}
% \State{Let $X=W=x'$ with probability $\alpha(x, p, x') = \alpha(x,x')$, otherwise stay.}

% \Else
% \State{Let $X=x'$ with probability $\alpha(x,x')$, otherwise stay.}
% \State{We keep $W=x$ at the original position.}
% \EndIf

% \end{algorithmic}
% \end{algorithm}

%\noindent
%Trivially, if we denote $X \sim P_X(x, \cdot)$ and $W \sim P_W(x, \cdot)$, then $P_X(x, \cdot) = \mathcal{L}_{\text{RWM}}(x, \cdot)$ and $P_W(x, \cdot) = \mathcal{L}_{\text{LGWM}}(x, \cdot)$ precisely. Coupling inequality states that
%\begin{equation}
%\label{eqn: TV inequality, RWM and LGWM}
%    \|\mathcal{L}_{\text{RWM}}(x,\cdot) - \mathcal{L}_{\text{LGWM}}(x, \cdot)\|_{TV}= \|P_X(x, \cdot) - P_W(x, \cdot)\|_{TV} \leq P(X \neq W)
%\end{equation}
%Note that the only possible case for $X \neq W$ is when $z > 0$ and $X$ does a transition to $x' > x$, the probability for this is computed as
% \begin{equation}
% \begin{aligned}
%     P(X \neq W) = P(z>0, X \text{ transition}) &= P(X \text{ transition} \mid z>0)P(z>0)\\
%     &= \frac{1}{2}\int_x^{\infty}\alpha(x,y)q(x,dy) = \frac{1}{2}\lim_{c\to \infty}\textcircled{1}
% \end{aligned}
% \end{equation}
% see the derivation of \textcircled{1} in \eqref{eqn: expression of textcircled 1}, along with \textcircled{1} $\to 0$ when $x \to \infty$ in \ref{appx subsec: computation of related terms}, we then instantly have
% \begin{equation}
% \label{eqn: TV coupling proof, RWM-LGWM, limit form}
%     \lim_{x\to \infty}P(X \neq W) = 0 
% \end{equation}
% Finally, by taking the limit on both sides of \eqref{eqn: TV inequality, RWM and LGWM}, and using \eqref{eqn: TV coupling proof, RWM-LGWM, limit form}, \eqref{eqn: TV convergence statement, LGWM and RWM} is therefore proved.
% \end{proof}


%\subsubsection{Multiple-step law of transition analysis}
%In the previous section, we have successfully proved that, for the light-tailed target distributions (and we use $\pi(x) = N(x;0,1)$ as a canonical example), LGWM performs rather in a same pattern as the classic RWM in terms of TV convergence for their first-step law of transition, written mathematically in terms of \eqref{eqn: TV convergence statement, LGWM and RWM}. 
%In this section we extend Proposition \ref{prop:1step_coupling} to the $n$-th step laws.

Recall that in general $\check{P}$ is called an $\varepsilon$-lazy version of $P$ if $\check{P}(x,\cdot) := (1-\varepsilon)P(x,\cdot) + \varepsilon \delta_x(\cdot)$ for all $x$.  The main result of this section is the following proposition, which shows that in the typical light-tailed scenario characterised by Assumption \ref{ass:strict_convexity} the random walk Metropolis behaves just as a $1/2$-lazy version of the guided walk Metropolis for large $|x|$.  

\begin{proposition}
\label{prop:nstep_coupling}
Let $\mathcal{L}(X_n)$ and $\mathcal{L}(W_n)$ denote the laws of the $n$th states of the random walk Metropolis chain $\{X_n\}_{n \geq 0}$ and the $1/2$-lazy guided walk Metropolis chain $\{(W_n,P_n)\}_{n\geq 0}$.  If $\pi$ satisfies Assumption \ref{ass:strict_convexity}, $X_0 = W_0 = x$ and $P_0 = -\operatorname{sgn}(x)$ then for any fixed $n \in \mathbb{N}$
\begin{equation*}
%\label{eqn: TV convergence statement, n-step, LGWM and RWM}
\lim_{|x|\to \infty} \|\mathcal{L}(X_n) - \mathcal{L}(W_n)\|_{TV} = 0.
\end{equation*}
If instead $X_0 = W_0 = x$ and $P_0 = \operatorname{sgn}(x)$, then
\begin{equation*}
\lim_{|x|\to \infty} \|\mathcal{L}(X_{n-1}) - \mathcal{L}(W_n)\|_{TV} = 0.
\end{equation*}
\end{proposition}

\begin{proof}[Proof of Proposition \ref{prop:nstep_coupling}]
We describe the case $x \to \infty$, and simply note that the $x \to -\infty$ case follows an identical argument.  We couple the random walk and the lazy guided walk proposals synchronously in the following way. Given the starting position $X_0 = x$ for the random walk and $(W_0, P_0) = (x,-1)$ for the lazy guided walk algorithms, draw $Z \sim N(0,1)$ and set the random walk proposal to be $x' = x - \epsilon Z$. The lazy guided walk proposal can be set as $(w',p') = (x - \epsilon \max(Z,0), -1)$, since $\mathbb{P}(\max(Z,0) = 0) = 1/2$, in which case the chain will be lazy and stay put. If $Z \geq 0$ then clearly $x' = w'$, so upon acceptance $X_1 = W_1 = x'$ and $P_1 = -1$.  If $Z < 0$ then there is a chance that the chains can separate if the random walk proposal is accepted, since the lazy guided walk chain will not move.  In Lemma \ref{lem:convergence_alpha_light_tails} below, however, we show that under Assumption \ref{ass:strict_convexity} the probability of accepting such a random walk proposal tends to $0$ as $x \to \infty$.  Similarly it is easy to show that when $Z > 0$ under the same assumption $\alpha(x,x') = 1$ whenever $Z \in [0, x/\epsilon]$ for sufficiently large $x$, and that $\mathbb{P}(Z < x/\epsilon | Z > 0) \to 1$ as $x \to \infty$.  It therefore holds that $\mathbb{P}(X_1 \neq W_1) \to 0$ as $x \to \infty$.  Iterating the argument to show that $\mathbb{P}(X_n \neq W_n) \to 0$ for any fixed and finite $n$ is tedious but straightforward, since a large enough $x$ can be chosen that $\mathbb{P}\left( x- \epsilon\sum_{i=1}^n \max(0, Z_i) > c\right)$ is arbitrarily close to 1 for any fixed $c < \infty$, where each $Z_i \sim N(0,1)$ represents the Gaussian draw used for the $i$th proposal of each algorithm. We omit the details for brevity.  Applying the coupling inequality then gives $\|\mathcal{L}(X_n) - \mathcal{L}(W_n)\|_{TV} \leq \mathbb{P}(X_n \neq W_n) \to 0$ as $x \to \infty$, which completes the proof in the case $P_0 = -\operatorname{sgn}(x)$.

In the case $P_0 = 1$ then by the same Lemma \ref{lem:convergence_alpha_light_tails} the initial guided walk proposal will be rejected with probability tending to 1 as $x \to \infty$. With probability tending to 1, therefore, $(W_1,P_1) = (x, -1)$. The same coupling as above can then be constructed to connect $X_{n-1}$ and $W_n$ in law as required.  This completes the proof.
\end{proof}


% If $pz \geq 0$ note also that $x' = w'$.  In this instance set $X_1 = W_1 = x'$ with probability $\alpha(x,x')$, otherwise let $X_1 = x$ and $(W_1, P_1) = (x, -p)$. If instead $pz < 0$ then $x' \neq w'$, but upon rejection $X_1 = W_1 = x$.  The next states for the position variable will therefore be equal if either $pz > 0$ and the proposals are both accepted, or in the case that both proposals are rejected.  Since in the former case $\alpha(x,x') = \alpha(x,w')$ then using a common $\mathrm{Unif}[0,1]$ draw to make the acceptance decision ensures that when $pz > 0$ then $X_1 = W_1$. To ensure that the laws become equal in general, we must show that if $pz \leq 0$ then the acceptance probability $\to 0$ as $x \to \infty$, and such probability is easily $\int_{x}^{\infty}\alpha(x,y)q(x, dy)/2$ (or $\int_{-\infty}^{x}\alpha(x,y)q(x, dy)/2$ if we consider the tail $(-\infty, x]$, with $x\to -\infty$), that converges to 0 by Lemma \ref{lem:convergence_alpha_light_tails}.  This establishes that $\mathbb{P}(X_1 \neq W_1) \to 0$ as $|x| \to \infty$.

% The extension 


% We construct a coupling of the two processes. Define the event
% $E_t := \{X_t \neq W_t\}$. The evolution for $t=1, \cdots, n$ is described as follows: when $X_{t-1} = W_{t-1} (= x_{t-1})$, we draw a single Gaussian noise $z \sim N(0,1)$ and form the random walk proposal $x'=x_{t-1} + \epsilon z$ and guided walk proposal $w' = x_{t-1} + \epsilon p_{t-1}|z|$ respectively; (i).  If $p_{t-1}z > 0$ then $x'=w'$, and in this case we compute $\alpha(x_{t-1}, x')$ and set $X_t = x'$ w.p. $\alpha(x_{t-1},x')$ and $(W_t, P_t) = (x', p_{t-1})$ w.p. INSERT, otherwise we keep $X_t = x_{t-1}$ and set $(W_t, P_t) = (x_{t-1}, -p_{t-1})$; (ii). if instead $p_{t-1}z \leq 0$, then $x'\neq w'$, and we still apply the RWM proposal that $X_{t} = x'$ with probability $\alpha(x_{t-1}, x')$ (otherwise it stays at $X_t = x_{t-1}$), while the LGWM chain remains at $(W_t, P_t) = (x_{t-1}, p_{t-1})$. On the other hand, when the chains are already different at the start of the step (i.e. $X_{t-1} \neq W_{t-1}$), then evolve them independently.

% % \begin{algorithm}[ht]
% % \caption{Sequential Coupling for RWM and LGWM, $x>0$}\label{alg: sequential coupling, RWM and LGWM}
% % \begin{algorithmic}[1]
% % \State{Given the initial state $x>0$, set $X_0=W_0 = x$, and $p_0 = -\operatorname{sgn}(x)$.}
% % \For{$t=1, \cdots, n$}
% %     \If{$X_{t-1} = W_{t-1} (=x_{t-1})$}
% %         \Comment{both $X_t$ and $W_t$ will generate with the same noise $z$}
% %         \State{Sample $z \sim N(0,1)$ independently.}
% %         \State{Compute $x'=X_{t-1} + \epsilon z$ and $w'=W_{t-1} + \epsilon p_{t-1} |z|$.}
% %         \If{$p_{t-1}z > 0$}
% %         \Comment{$p_{t-1}z>0$ ensures $x'=w'$ always}
% %             \State{Let $X_t = W_t = x'$ with probability $\alpha(x_{t-1}, p_{t-1}, x') = \alpha(x_{t-1}, x')$, otherwise stay.}
% %             \State{Assert $p_t=p_{t-1}$ with probability  $\alpha(x_{t-1}, p_{t-1}, x') = \alpha(x_{t-1}, x')$, otherwise flip.}
% %         \Else
% %             \State{Let $X_t = x'$ with probability $\alpha(x_{t-1}, x')$, otherwise stay.}
% %             \State{We keep $W_t = W_{t-1}=x_{t-1}$ at the previous position, also keep $p_t=p_{t-1}$.}
% %         \EndIf
% %     \Else
% %         \State{Generate $X_t \sim \mathcal{L}_{\text{RWM}}(X_{t-1}, \cdot)$ and $W_t \sim \mathcal{L}_{\text{LGWM}}(W_{t-1}, \cdot)$ (with $p_{t-1}$) independently.}    
% %     \EndIf
% % \EndFor

% % \end{algorithmic}
% % \end{algorithm}

% Note that given $E_{t-1}^c$ with known $p_{t-1}$, such disagreement at $t$ can only occur at $p_{t-1}z\leq 0$ with $X_t=x'$ transitioned but $W_t = W_{t-1} = X_{t-1}$ stayed. We mathematically have,
% \begin{equation*}
%     P(p_{t-1}z \leq 0) = P(z>0)\mathbf{1}(p_{t-1} = -1) + P(z\leq0)\mathbf{1}(p_{t-1} = 1) = \frac{1}{2}
% \end{equation*}
% which reflects the lazy ratio 1/2. Therefore, when we treat $x_{t-1}$ and $p_{t-1}$ as deterministic values
% \begin{equation}
% \label{eqn: P(E_t|E_{t-1}^C)}
% \begin{aligned}
%     &P(E_t \mid E_{t-1}^c) = P(p_{t-1}z \leq 0, X_t \text{ transition}) = P(X_t \text{ transition} \mid p_{t-1}z \leq 0)P(p_{t-1}z \leq 0)\\
%     &= \frac{1}{2}\underbrace{\int_{x_{t-1}}^\infty \alpha(x_{t-1},y)q(x_{t-1},dy)\mathbf{1}(p_{t-1} = -1)}_{\text{tail $[x_{t-1}, \infty)$ as $z>0 \implies y=x_{t-1}+\epsilon z > x_{t-1}$}} + \frac{1}{2}\underbrace{\int_{-\infty}^{x_{t-1}}\alpha(x_{t-1},y)q(x_{t-1},dy)\mathbf{1}(p_{t-1} = 1)}_{\text{tail $(-\infty, x_{t-1}]$ as $z\leq 0 \implies y=x_{t-1}+\epsilon z \leq x_{t-1}$}}
% \end{aligned}
% \end{equation}
% We first discuss the case when $p_{t-1} = -1$ for the tail region $[x_{t-1}, \infty)$ with $x_{t-1} > 0$. 
% \begin{equation*}
% \begin{aligned}
%     P(E_t \mid E_{t-1}^c) &= \frac{1}{2}\int_{x_{t-1}}^\infty \alpha(x_{t-1},y)q(x_{t-1},dy), \quad \text{and } y=x_{t-1}+\epsilon z\\
%     &= \frac{1}{2}\int_{0}^{\infty}\alpha(x_{t-1}, x_{t-1}+\epsilon z)\phi(z)dz
% \end{aligned}
% \end{equation*}
% where $\phi(z) = (2\pi)^{-1/2}\exp(-z^2/2)$ is the density function for $N(0,1)$. As shown in the proof of \ref{lem:convergence_alpha_light_tails} in the Appendix \ref{app:proofs_1D_RWM_GWM}, applying the mean value theorem gives us
% \begin{equation*}
%     U(x_{t-1}+\epsilon z) - U(x_{t-1}) \geq \epsilon z \inf_{t \in [x_{t-1}, x_{t-1}+\epsilon z]} U'(t) = \epsilon z U'(x_{t-1})
% \end{equation*}
% This subsequently indicates
% \begin{equation*}
%     \alpha(x_{t-1}, x_{t-1}+\epsilon z) \leq \min\left\{1, \exp\left(-\epsilon z U'(x_{t-1})\right)\right\}
% \end{equation*}
% Given that we are considering the tail region $[x_{t-1}, \infty)$ with the setting that $x_{t-1} > 0$, since $U'$ is continuous and monotonic increasing, with $U'(x) \to \pm \infty$ when $x\to \pm \infty$ respectively, we can definitely find a $K>0$ such that $U'(K) > 0$ and $U'(-K) \leq 0$ (GRAPHICAL ILLUSTRATION?). Therefore, if $x_{t-1} \geq K > 0 \implies U'(x_{t-1}) \geq U'(K) > 0$, then for the tail region $[x_{t-1}, \infty)$ we shall reduce the acceptance probability as $\alpha(x_{t-1}, x_{t-1}+\epsilon z) \leq \exp\left(-\epsilon z U'(x_{t-1})\right)$, and then, clearly that $\phi(z) \leq (2\pi)^{-1/2}$, we shall have the following inequality that
% \begin{equation*}
% \begin{aligned}
%     P(E_t \mid E_{t-1}^c) &= \frac{1}{2}\int_{0}^{\infty}\alpha(x_{t-1}, x_{t-1}+\epsilon z)\phi(z)dz\\
%     &\leq \frac{1}{2\sqrt{2\pi}}\frac{1}{\epsilon U'(x_{t-1})}
% \end{aligned}
% \end{equation*}
% Using similar logic, for $p_{t-1}=1$ we can express the acceptance probability for the tail region $(-\infty, x_{t-1}]$ under $x_{t-1} \leq -K \leq 0 \implies U'(x_{t-1}) \leq U'(-K) \leq 0$ in terms of $\alpha(x_{t-1}, x_{t-1}-\epsilon z) \leq \exp\left(\epsilon z U'(x_{t-1})\right)$, and hence
% \begin{equation*}
% \begin{aligned}
%     P(E_t \mid E_{t-1}^c) &= \frac{1}{2}\int_{-\infty}^{0}\alpha(x_{t-1}, x_{t-1}+\epsilon z)\phi(z)dz = \frac{1}{2}\int_0^{\infty} \alpha(x_{t-1}, x_{t-1}-\epsilon u)\phi(u)du\\
%     &\leq -\frac{1}{2\sqrt{2\pi}}\frac{1}{\epsilon U'(x_{t-1})}
% \end{aligned}
% \end{equation*}
% Combining both cases that, for the positive tail region $[x_{t-1}, \infty)$ under $x_{t-1}>0$ or the negative tail region $(-\infty, x_{t-1}]$ under $x_{t-1}\leq 0$, when we denote some $K>0$ that $U'(K) > 0$ and $U'(-K)\leq 0$, then for $|x_{t-1}|> K$ we have
% \begin{equation}
% \label{eqn: upper bound for 1-step disagreement, multistep RWM and LGWM}
%     P(E_t \mid E_{t-1}^c)\leq \frac{C}{\epsilon |U'(x_{t-1})|} \leq \frac{C}{\epsilon K^*}, \quad C = \frac{1}{2\sqrt{2\pi}}
% \end{equation}
% where $K^*:= \min(U'(K), |U'(-K)|)$. Note that \eqref{eqn: upper bound for 1-step disagreement, multistep RWM and LGWM} is only true if we are at the correct tail, i.e. $x_{t-1}p_{t-1} \leq 0 \iff p_{t-1} = -\operatorname{sgn}(x_{t-1})$.
% % Moreover, clearly that $\phi(z) \leq (2\pi)^{-1/2}$, then we shall have the following inequality that
% % \begin{equation*}
% % \begin{aligned}
% %     P(E_t \mid E_{t-1}^c) &= \frac{1}{2}\int_{0}^{\infty}\alpha(x_{t-1}, x_{t-1}+\epsilon z)\phi(z)dz\\
% %     &\leq \frac{1}{2\sqrt{2\pi}}\frac{1}{\epsilon U'(x_{t-1})}\mathbf{1}(U'(x_{t-1})>0) + \frac{1}{4}\mathbf{1}(U'(x_{t-1}) \leq 0)
% % \end{aligned}
% % \end{equation*}
% % By similar logic, we can extend such an upper bound into the case where $p_{t-1} = 1$. Note that for $z > 0$ with strict monotonicity of $U'$ we have
% % \begin{equation*}
% %     U(x_{t-1}) - U(x_{t-1}-\epsilon z) \leq \epsilon z \sup_{t\in [x_{t-1}-\epsilon z, x_{t-1}]} U'(t) = \epsilon z U'(x_{t-1})
% % \end{equation*}
% % hence $\alpha(x_{t-1}, x_{t-1}-\epsilon z) \leq \min\left\{1, \exp\left(\epsilon z U'(x_{t-1})\right)\right\}$, and for $p_{t-1}=-1$ we have
% % \begin{equation*}
% % \begin{aligned}
% %     P(E_t \mid E_{t-1}^c) &= \frac{1}{2}\int_{-\infty}^{0}\alpha(x_{t-1}, x_{t-1}+\epsilon z)\phi(z)dz = \frac{1}{2}\int_0^{\infty} \alpha(x_{t-1}, x_{t-1}-\epsilon u)\phi(u)du\\
% %     &\leq -\frac{1}{2\sqrt{2\pi}}\frac{1}{\epsilon U'(x_{t-1})}\mathbf{1}(U'(x_{t-1}) \leq 0) + \frac{1}{4}\mathbf{1}(U'(x_{t-1}) > 0)
% % \end{aligned}
% % \end{equation*}
% % Combining both cases then gives us
% % \begin{equation}
% % \label{eqn: upper bound for 1-step disagreement, multistep RWM and LGWM}
% %     P(E_t \mid E_{t-1}^c)\leq \begin{cases}
% %         \frac{1}{4} &\hbox{if $U'(x_{t-1})p_{t-1} \leq 0$}\\
% %         \frac{C}{\epsilon |U'(x_{t-1})|} &\hbox{otherwise}
% %     \end{cases}
% % \end{equation}
% % with $C = 1/2\sqrt{2\pi}$. 

% % we can always find one $x$ such that $U'(x) > 0$, and the statement also holds for $U'(x) < 0$, based on the fact that $U'$ is continuous and $U'(x) \to \pm\infty$ when $x\to \pm \infty$ respectively. This allows us to define $R$.

% Let's now consider the full history up to time $t$, with all intermediate states $X_1, \cdots, X_{n}$ being as random variables. Our interest now would be investigating that, given a very large (in magnitude) initial position $|x| \to \infty$, whether we still have very high probability that for all $t = 1, \cdots, n$, $|X_t|$ being large enough. To do this, define a region $R$ that all intermediate states of $X_t$ lie within this region
% \begin{equation}
% \label{eqn: region R}
%     R = \{|X_t| > K: \forall t=0, 1,\cdots, n\}
% \end{equation}
% for some large enough $K$ chosen $U'(K) > 0$ and $U'(-K) \leq 0$. Within this $R$, we claim the following lemma.

% Lemma \ref{lemma: stability of momentum in the tail region} simply states that once we start from the tail $|x|\to \infty$, within $R$ we will have a very high probability for $P_{t} = -\operatorname{sgn}(X_t)$, for all $t\in \mathbb{N}$ given that all previous intermediate states are equal. Put this into \eqref{eqn: P(E_t|E_{t-1}^C)}, it simply means we will discuss the correct tail region with high probability and then by \eqref{eqn: upper bound for 1-step disagreement, multistep RWM and LGWM} we immediately have
% \begin{equation}
% \label{eqn: upper bound with full history within R, RWM and LGWM}
%     P(E_t \mid E_{t-1}^c \cap E_{t-2}^c \cap \cdots E_1^c \cap R) \leq \frac{C}{\epsilon K^*}, \quad K^*= \min(U'(K), |U'(-K)|)
% \end{equation}


% % \begin{equation}
% % \label{eqn: upper bound with full history, multistep RWM and LGWM}
% %     P(E_t \mid E_{t-1}^c \cap E_{t-2}^c \cap \cdots E_1^c \cap R) \leq \sup_{X_{t-1} \in \mathbb{R}}P(E_t\mid X_{t-1} = W_{t-1})\leq \sup_{X_{t-1}\in \mathbb{R}}\frac{C}{\epsilon |U'(X_{t-1})|}
% % \end{equation}


% % With \eqref{eqn: upper bound for 1-step disagreement, multistep RWM and LGWM}, if we know the exact value of $x_{t-1}$ that results $X_{t-1} = W_{t-1} = x_{t-1}$, then we shall have
% % \begin{equation*}
% %     P(E_t \mid E_{t-1}^c \cap E_{t-2}^c \cap \cdots E_1^c) = P(E_t\mid X_{t-1} = W_{t-1}=x_{t-1})\leq \frac{C}{\epsilon |U'(x_{t-1})|}
% % \end{equation*}
% % However, through our coupling such intermediate state $X_{t-1} = x_{t-1} \in \mathbb{R}$ is just random. Such randomness would result our upper bound for the full history written as
% % \begin{equation}
% % \label{eqn: upper bound with full history, multistep RWM and LGWM}
% %     P(E_t \mid E_{t-1}^c \cap E_{t-2}^c \cap \cdots E_1^c) \leq \sup_{X_{t-1} \in \mathbb{R}}P(E_t\mid X_{t-1} = W_{t-1})\leq \sup_{X_{t-1}\in \mathbb{R}}\frac{C}{\epsilon |U'(X_{t-1})|}
% % \end{equation}
% % It is obvious that if every intermediate state has very large magnitude (i.e. $|x_{t}|\to \infty$ for $t=1, \cdots, n$), then $|U'(x_t)|\to \infty$ (see the discussion of this in the proof of Lemma \ref{lem:convergence_alpha_light_tails} in the Appendix \ref{app:proofs_1D_RWM_GWM}) and we shall have $P(E_t \mid E_{t-1}^c)\to 0$. 

% % Our interest now would be investigating that, given a very large (in magnitude) initial position $|x| \to \infty$, whether we still have `very high probability' that for all $t = 1, \cdots, n$, $|X_t|$ being `large enough'. To do this, define a region $R$ that all intermediate states of $X_t$ lie within this region
% % \begin{equation}
% % \label{eqn: region R}
% %     R = \{|X_t| > K: \forall t=1,\cdots, n\}
% % \end{equation}
% % for some large enough $K$ chosen $U'(K) > 0$ and $U'(-K) \leq 0$. Within this $R$, we can relax \eqref{eqn: upper bound for 1-step disagreement, multistep RWM and LGWM} and \eqref{eqn: upper bound with full history, multistep RWM and LGWM} with the common upper bound that
% % \begin{align}
% % \label{eqn: upper bound for 1-step disagreement within R, RWM and LGWM}
% %     &P(E_t \mid E_{t-1}^c \cap R) \leq \frac{C}{\epsilon K}\\
% % \label{eqn: upper bound with full history within R, RWM and LGWM}
% %     &P(E_t \mid E_{t-1}^c \cap E_{t-2}^c \cap \cdots E_1^c \cap R) \leq \frac{C}{\epsilon K}
% % \end{align}

% \noindent
% With \eqref{eqn: upper bound with full history within R, RWM and LGWM} we can then represent the first time disagreement probability 
% \begin{equation}
% \label{eqn: multistep proof inequality 1}
% \begin{aligned}
%     P(\exists t \leq n:E_t \mid R) &= P\left(\bigcup_{t=1}^n (E_t \cap E_{t-1}^c \cap \cdots \cap E_1^c) \mid R\right)\\
%     &\leq \sum_{t=1}^n P(E_t \cap E_{t-1}^c \cap \cdots \cap E_1^c \mid R)\\
%     &\leq \sum_{t=1}^n P(E_t \mid E_{t-1}^c \cap \cdots \cap E_1^c \cap R) \leq \frac{nC}{\epsilon K^*}
% \end{aligned}
% \end{equation}
% where the third inequality is due to $P(A \cap B) = P(A \mid B)P(B) \leq P(A \mid B)$ for $P(B) \leq 1$, and we apply \eqref{eqn: upper bound with full history within R, RWM and LGWM} for the last inequality. 

% Meanwhile, we compute the $P(R^c)$, which is the probability for existing intermediate states that move into a compact, closed ball centered 0 with radius $K$. By the same logic as \eqref{eqn: multistep proof inequality 1}, with given initial position $X_0=x$ (so $P(X_0=x) = 1$) we can write
% \begin{equation*}
% \begin{aligned}
%     P(R^c) &= P(\exists t\leq n: |X_t|\leq K) = P\left[\left(\bigcup_{t=1}^n |X_t| \leq K \right) \mid X_0=x\right]P(X_0 = x)\\
%     &\leq \sum_{t=1}^n P(|X_t|\leq K \mid X_0=x)
% \end{aligned}
% \end{equation*}
% We now claim the following lemma.

% Therefore, with Lemma \ref{lemma:out_of_R_zero_prob}, we can then write
% \begin{equation}
% \label{eqn: multistep proof limit P(R^c)}
%     \lim_{x\to \pm\infty}P(R^c) = 0 \iff \forall \kappa > 0, \exists \delta = \delta(\kappa) >0, \forall |x| > \delta, P(R^c) < \kappa
% \end{equation}
% Finally, by \eqref{eqn: multistep proof inequality 1} we have
% \begin{equation}
% \begin{aligned}
%     P(E_n) &= P(E_n \cap R) + P(E_n \cap R^c)\\
%     &\leq P(E_n \mid R) + P(R^c)\\
%     &\leq P(\exists t \leq n : E_t \mid R) + P(R^c) \leq \frac{nC}{\epsilon K^*} + P(R^c)
% \end{aligned}
% \end{equation}
% easily as $\{E_n \mid R\} \subseteq \{\exists t\leq n: E_t \mid R\}$. For any $\kappa > 0$, we can always choose large enough $K = K(\kappa)$ such that with given $n, C, \epsilon$, we have $nC/\epsilon K^* < \kappa/2$, meanwhile, by \eqref{eqn: multistep proof limit P(R^c)} we know there must exist $\delta = \delta(\kappa)$ such that for all $|x| > \delta$ we have $P(R^c) < \kappa/2$. Together, it can be indeed written as
% \begin{equation}
%     \forall \kappa > 0, \exists K = K(\kappa), \delta = \delta(\kappa), \forall |x| > \delta, P(E_n) < \kappa/2 + \kappa/2 = \kappa \iff \lim_{x\to \pm\infty}P(E_n)=0
% \end{equation}
% Recall that coupling inequality allows us to write
% \begin{equation*}
%     \|\mathcal{L}_{\text{RWM}}^n(x, \cdot) - \mathcal{L}_{\text{LGWM}}^n(x, \cdot)\|_{TV} \leq P(X_n\neq W_n) = P(E_n) \to 0, \quad \text{when }x\to \pm \infty
% \end{equation*}
% This completes the proof for Proposition \ref{prop:nstep_coupling}.
% \end{proof}



% \begin{lemma}[Stability of momentum in the tail region]
% \label{lemma: stability of momentum in the tail region}
%     Let $P_0 = -\operatorname{sgn}(x)$ as the deterministic initial momentum, and define the event $A_t = \{P_t = -\operatorname{sgn}(W_t)\}$, also region $R$ as \eqref{eqn: region R}. For a general light-tailed target and any fixed number of steps $n$, as $|x| \to \infty$, we have
% \begin{equation}
%     P\left(A_t \mid E_{t-1}^c \cap E_{t-2}^c \cap \cdots \cap E_1^c \cap R\right) \to 1, \quad \forall t \in \mathbb{N}
% \end{equation}
% \end{lemma}

% \begin{proof}
%     See Appendix \ref{app:proofs_1D_RWM_GWM}.
% \end{proof}

% \begin{lemma}
% \label{lemma:out_of_R_zero_prob}
%     With the region $R$ with boundary $K$ defined in \eqref{eqn: region R}, also the target being as general light-tailed, we have $\lim_{x\to \pm \infty}P(|X_t|\leq K \mid X_0=x) = 0$.
% \end{lemma}

% \begin{proof}
%     See Appendix \ref{app:proofs_1D_RWM_GWM}.
% \end{proof}


\section{Discussion}

We have shown that non-reversible Metropolis--Hastings algorithms can sometimes produce much faster mixing chains than reversible counterparts, but that in other cases they behave similarly.  In particular, a random walk Metropolis will not always exhibit `random walk' behaviour; this depends crucially on the form of $\pi$, specifically the tails, which determine how much the Metropolis--Hastings chain $P$ is altered compared to the random walk proposal $Q$.  Non-reversible chains appear to be most advantageous when $\pi$ has heavy tails or more generally becomes flat in some region or direction.

Extensions of the guided walk Metropolis to $\mathbb{R}^d$ have been constructed but tend to be inferior to other (reversible) approaches \citep{gustafson2004value}.  Continuous-time processes that are similar in spirit have, however, shown more success in high dimensions \citep{bierkens2019zig}.  Another super-diffusive process is a random walk with heavy-tailed increments (see e.g. \citep{metzler2000random,klafter1987stochastic}).  Heavy-tailed increments can often be induced in sampling algorithms via randomised step-sizes, which are commonly employed and known to have benefits \cite{grazzi2026randomized}.  Using a heavy-tailed proposal in the random walk Metropolis can also increase the polynomial rate of convergence when $\pi$ has polynomial tails \cite{jarner2007convergence}.  We suspect that a similar improvement in the rate of a modified guided walk algorithm in which $Z$ is heavy-tailed can be shown, but leave the details for future work.  
%=========================================

\subsection*{Acknowledgments}

The authors would like to thank the Isaac Newton Institute for Mathematical Sciences, Cambridge, for support and hospitality during the programme \textit{Stochastic systems for anomalous diffusion}, where work on this paper was undertaken. This work was supported by EPSRC grant EP/Z000580/1.

\clearpage
%% reset page counter and start appendix pages numbering
\pagenumbering{arabic}
\renewcommand*{\thepage}{References | Page \arabic{page}}

%%References part of the main text
% References: modify the file refs.bib
\bibliographystyle{abbrvnat}
\bibliography{refs}

\clearpage
%% reset page counter and start appendix pages numbering
\pagenumbering{arabic}
\renewcommand*{\thepage}{A.\arabic{page}}
\renewcommand{\thesection}{\Alph{section}}
%% Supplementary Material goes here
\appendix


%%%%%%%%%%%%%% Section 2 %%%%%%%%%%%%%%%%%%%%%%%
\section{Proofs of Section \ref{sec:MH_high_accept}}
\label{app:proofs_MH_high_accept}

\begin{proof}[Proof of Proposition \ref{prop: MH is a T-chain}]
(i) \textit{$T$-chain:} First note that choosing $\nu = \delta_1$ gives 
\begin{equation*}
    P(x,A) \ge\int_A \alpha(x,y)q(x,y)dy =: T(x,A)
\end{equation*}
for any $A \in \mathcal{B}(\mathbb{R}^d)$. Define the function
\begin{align}
    T_y(\cdot) &:= \alpha(\cdot,y)q(\cdot,y) = \min\left\{q(\cdot,y),\frac{\pi(y)q(y,\cdot)}{\pi(\cdot)}\right\},
\end{align}
meaning $y$ is treated as a constant and $T_y(x)$ is viewed as a function of $x$ only.  By assumption $\pi(x)$ is upper semicontinuous and positive, meaning $1/\pi(x)$ is lower semicontinuous, as is the function $q(y,\cdot)$. Therefore for any sequence $x_n\to x$
$$
\liminf_{n\to\infty} \frac{q(y,x_n)}{\pi(x_n)}\ge \left( \liminf_{n\to\infty} q(y,x_n)\right) \cdot \left( \liminf_{n\to\infty} \frac{1}{\pi(x_n)}\right)\ge \frac{q(y,x)}{\pi(x)},
$$
meaning $\pi(y) q(y,\cdot)/\pi(\cdot)$ is lower semicontinuous.  The minimum of two lower semicontinuous functions is also lower semicontinuous, since for each $c \in \mathbb{R}$
$$
\{x\in \mathcal{X}:\min\{g(x),h(x)\}\le c\} =\{x\in \mathcal{X}:g(x)\le c\} \cup \{x\in \mathcal{X}:h(x)\le c\},
$$
which is closed if $g$ and $h$ are both lower semicontinuous. Combining gives that $T_y(x)$ is lower semicontinuous in $x$.  Upon noting that $T(x,A) = \int_A T_y(x) dy$, then for any sequence $x_n\to x$
\begin{align}
    \liminf_{n\to\infty} T(x_n,A) = \liminf_{n\to\infty}  \int T_y(x_n)\mathbb{I}_A(y)dy
    \ge   \int \liminf_{n\to\infty} T_y(x_n)\mathbb{I}_A(y)dy
    \ge \int_A T_y(x)dy  = T(x,A)
\end{align}
using Fatou's lemma and lower semicontinuity of $T_y(x)$, which shows that $T(\cdot, A)$ is lower semicontinuous for any $A \in \mathcal{B}(\mathbb{R}^d)$. Clearly $T(x,\mathbb{R}^d) \geq \int_{\|y-x\|\leq \delta_q}q(x,y)\alpha(x,y)dy > 0$ for any $x$ by Assumption \ref{ass:pi_q_regularities} and therefore $\{X_n\}_{n\geq 0}$ is a $T$-chain.

\textit{$\pi$-irreducibility:} 
By assumption if $\|x-y\|\le \delta_q$, then $q(x,y)>0$. This together with the assumptions on $\pi(x)$ implies
\begin{equation} \label{eq:q_alpha_lb}
q(x,y)\alpha(x,y) = \min\left(q(x,y),\frac{\pi(y)q(y,x)}{\pi(x)}\right) > 0.
\end{equation}
%For any $x, z \in \mathbb{R}^d$ there exists a path of finite length $n$ and points $x_0, x_1, \dots, x_n$ with $x_0 = x, x_n = z$ such that $\|x_i - x_{i-1}\| < \delta_q$.  Combining with \eqref{eq:q_alpha_lb} gives the sub-transition density $p^n(x,z)$ associated with generating $z$ from $x$ in $n$ steps satisfies
%\[
%p^n(x, z) \geq \int \dots \int \prod_{i=1}^n q(x_{i-1}, x_i)\alpha(x_{i-1},x_i) dx_1 \dots dx_{n-1} > 0
%\]
%Since $z$ is arbitrary then $\pi(A) > 0$ implies $A \in \mathcal{B}(\mathbb{R}^d)$ has positive Lebesgue measure, it follows that $P^n(x, A) \ge  \int_A p^n(x,z)dz \ge > 0$. The chain is therefore $\pi$-irreducible.
Consider any set $A \in \mathcal{B}(\mathbb{R}^d)$ such that $\pi(A) > 0$. Since $\pi$ has a density that is positive on all of $\mathbb{R}^d$, $A$ must have positive Lebesgue measure. The inner regularity property of the latter implies that $A$ contains a compact subset $K \subseteq A$ also of positive Lebesgue measure. For a fixed $x$ and any $z \in K$, there exists a path of finite length $n$ and points $x = x_0, x_1, \dots, x_n = z$ such that $\|x_i - x_{i-1}\| < \delta_q$. Since $K$ is bounded, we can choose an $n$ large enough to reach any $z \in K$ from $x$. The sub-transition density $p^n(x,z)$ associated with generating $z$ from $x$ in $n$ steps satisfies
\[
p^n(x, z) \geq \int \dots \int \prod_{i=1}^n q(x_{i-1}, x_i)\alpha(x_{i-1},x_i) dx_1 \dots dx_{n-1} > 0
\]
for all $z \in K$. It then follows that the transition probability to the set $A$ satisfies
\[
P^n(x, A) \ge P^n(x,K) \geq \int_K p^n(x,z)dz > 0
\]
from which it follows that the chain is $\pi$-irreducible.

\textit{Aperiodicity:}
To show aperiodicity, it suffices to prove that for all $x\in \mathbb{R}^d$, there exists $A \in \mathcal{B}(\mathbb{R}^d)$ such that both $P(x,A)>0$ and $P^2(x,A)>0$.  Let $B_x =\{y:\|y-x\|<\delta_q\}$, and note that $P(x,B_x) >0$ by Assumption \ref{ass:pi_q_regularities}.  For $P^2(x,B_x)$, we have
\begin{equation}
    P^2(x,B_x) = \int P(x,dy)P(y,B_x) \ge \int_{B_x} q(x,y)\alpha(x,y)P(y,B_x)dy.
\end{equation}
For $y\in B_x$, define $B_y  = \{z:\|z-y\|<\delta_q\}$. Since $\|y-x\| < \delta_q$, the intersection $B_x\cap B_y$ is nonempty. Therefore,
$$
P(y,B_x)\ge P(y, B_x\cap B_y) \ge \int_{B_x\cap B_y} q(y,z)\alpha(y,z)dz>0.
$$
Since the integrand $p(x,y)P(y,B_x)$ is the product of two strictly positive terms for all $y \in B_x$, we conclude that $P^2(x,B_x) > 0$, showing that the chain is aperiodic.

\textit{Smallness of compacts:}
By Theorem 5.5.7 and Theorem 6.2.5 in \cite{meyn2012markov}, if a chain is $\pi$-irreducible and aperiodic, then every petite set is small, and if the chain is also a $T$-chain, then every compact set is petite. Therefore, for a  $\pi$-irreducible and aperiodic $T$-chain, every compact set is small.

(ii) By part (i) and Assumption \ref{ass:pi_q_regularities} the conditions of Proposition 3.1 in \cite{roberts1996geometric} are satisfied, which implies the result.
\end{proof}


\begin{proof}[Proof of Proposition \ref{prop: counterexample}]
First note that
\begin{align*}
   \lim_{|x|\to\infty} \int\alpha(x,y) Q(x,dy)
    % & = \lim_{|x|\to\infty}\int\min\left\{Q(x,dy),\left[\left(1-\frac{1}{2+|y|}\right)\frac{1}{\sqrt{2\pi}}\exp\left(\frac{x^2}{2}+\frac{xy}{2}-\frac{9y^2}{8}\right) + \frac{\exp(x^2-y^2)}{2+|y|}\mathbf{1}_x(\exp(y^2),\exp(y^2)+1)\right]dy\right\}\\
    \ge \lim_{|x|\to\infty}\int \min\left\{Q(x,dy),\frac{1}{2\sqrt{2\pi}}\exp\left(\frac{x^2}{2}+\frac{xy}{2}-\frac{9y^2}{8}\right)dy\right\}= 1,
\end{align*}
To verify that $Q$ is not geometrically ergodic, suppose for contradiction that it is. By Proposition \ref{prop: MH is a T-chain}, there exists $V\ge1$ such that 
\[
\limsup_{|x|\to\infty}\frac{QV(x)}{V(x)} = \limsup_{|x|\to\infty} \left[\frac{(1-\varepsilon(x))Q_1V(x)}{V(x)} + \varepsilon(x)\frac{Q_2V(x)}{V(x)}\right]<1
\]
Since $\varepsilon(x)Q_2V(x)/V(x)\ge 0$, this inequality implies $\limsup_{|x|\to\infty}Q_1V(x)/V(x)<1$.  This implies further that $Q_1V(x) \leq \rho V(x)$ for sufficiently large $|x|$ and some $\rho < 1$.  In Lemma \ref{lem:ar1_vgrowth} we show that any such $V(x) \geq c|x|^\eta$ for all sufficiently large $|x|$. 
% Indeed, if $V(x) = c\log(|x|+e)$,where $c\ge 1$ then
% \begin{align}
%     \lim_{|x|\to\infty} \frac{Q_1V(x)}{V(x)}=\lim_{|x|\to\infty}\frac{\mathbb{E}_Z[\log (|x/2+Z|+e)]}{\log (|x|+e)} ,\quad Z\sim N(0,1)
% \end{align}
% Then by positivity of $V$ and letting $\rho$ denote the standard Normal density and $C_x = \{z,|z|\in[|x|/4,3|x|/4]\}$
% \begin{align*}
%     \mathbb{E}_Z\left[\log\left(\left|\frac{x}{2}+Z\right|+e\right)\right] &\geq  
%     \int_{\mathbb{R} \setminus C_x}\log\left(\left|\frac{x}{2}+z\right|+e\right)\rho(z)dz
%     \ge 
%     \log\left|\frac{x}{4}\right|\left(1-\int_{C_x}\rho(z) dz\right)
% \end{align*}
% The last inequality follows since for all $z \notin C_x$ it holds that $|x/2+z|>|x/4| $. It is clear that $\lim_{|x|\to\infty}\int_{C_x}\rho(z)dz = 0$, meaning that
% \begin{equation*}
%     \lim_{|x|\to\infty} \frac{Q_1V(x)}{V(x)}\ge \lim_{|x|\to\infty} \frac{\log(|x|/4)}{\log (2|x|)}\left(1-\int_{C_x}\rho(z)dz \right) = \lim_{|x|\to\infty} \left[1- \frac{\log 8}{\log 2|x|}\right] =1
% \end{equation*}
% Therefore, any suitable $V(x)$ must grow asymptotically faster than $\log|x|$. 
Note also that
\begin{align*}
    \limsup_{|x|\to\infty}\frac{QV(x)}{V(x)} 
    % &= \lim_{|x|\to\infty} \left[\frac{(1-\varepsilon(x))Q_1V(x)}{V(x)} + \varepsilon(x)\frac{Q_2V(x)}{V(x)}\right]\nonumber\\
    &\ge \limsup_{|x|\to\infty} \left[\varepsilon(x)\frac{Q_2V(x)}{V(x)}\right].
    %\ge \lim_{|x|\to\infty} \left[\frac{1}{2+|x|}\frac{x^2}{\log 2|x|}\right] =\infty,
\end{align*}
Using this and the condition $QV(x) \leq \rho V(x)$ for large enough $|x|$ implies
\[
\varepsilon(x)Q_2V(x) \leq QV(x) \leq \rho V(x),
\]
meaning we have the upper bound $Q_2V(x) \leq \rho(2+|x|)V(x)$.  
% Since $Q_2V \geq 0$, a similar argument gives $(1-\varepsilon(x))Q_1V(x) \leq \rho V(x)$, and since $\varepsilon(x) \to 0$ this means $Q_1V(x) \leq \rho' V(x)$ for large $|x|$ and some $\rho'<1$. Lemma \ref{lem:ar1_vgrowth} can then be applied giving
% $V(x) \geq c|x|^\eta$ for all large $|x|$.  
Applying the lower bound at any $y \in (\exp(x^2), \exp(x^2)+1)$ gives $V(y) \geq c\exp(\eta x^2)$. Integrating over the unit interval gives
\[
Q_2V(x) = \int_{\exp(x^2)}^{\exp(x^2)+1} V(y)\,dy \geq c\exp(\eta x^2).
\]
Combining with the upper bound leads to the inequality
\[
\rho(2+|x|)V(x) \geq c\exp(\eta x^2) \implies V(x) \geq \frac{c}{\rho(2+|x|)}\exp(\eta x^2).
\]
Now, applying this new lower bound at the point $y = \exp(x^2)$ gives
\[
V(\exp(x^2)) \geq \frac{c}{\rho(2+\exp(x^2))}\exp\bigl(\eta e^{2x^2}\bigr) \geq \exp\bigl(\eta' e^{2x^2}\bigr)
\]
for any $\eta' < \eta$ and large enough $|x|$.  This provides a lower bound for $Q_2V(x)$, leading to
\[
\exp(\eta' e^{2x^2}) \leq Q_2V(x) \leq \rho(2+|x|)V(x) \implies V(x) \geq \exp(\eta'' e^{2x^2})
\]
for $\eta''< \eta'$ in the same $|x|$ limit.  For any $V$ that satisfies the above, however, letting $Z \sim N(0,1)$ it holds that
\[
Q_1V(x) 
= 
\mathbb{E}[V(x/2+Z)] 
\geq 
\mathbb{E}\left[\exp\left(\eta''  e^{2(x/2+Z)^2}\right)\right] 
= 
\frac{1}{\sqrt{2\pi}}\int \exp\left(\eta'' e^{2(x/2+z)^2} - \frac{z^2}{2}\right)dz = \infty.
\]
This directly contradicts the requirement $Q_1V(x) \leq \rho V(x) < \infty$, from which we conclude that $Q$ cannot be geometrically ergodic.

To see that $P(x,dy)$ is geometrically ergodic, choosing $V(x) = \exp(x^2/4) $, note that
\begin{align*}
\limsup_{|x|\to\infty} \frac{PV(x)}{V(x)} & =\limsup_{|x|\to\infty}\left[\frac{\int V(y)\alpha(x,y)Q(x,dy) }{V(x)}+ \int(1-\alpha(x,y)) Q(x,dy) \right]\\
 &= \limsup_{|x|\to\infty}\frac{\int\alpha(x,y)(1-\varepsilon(x))V(y)Q_1(x,dy)  + \int_{\exp(x^2)}^{\exp(x^2)+1}\alpha(x,y)\varepsilon(x) V(y)}{V(x)}.
\end{align*}
Upon noting that $1-\varepsilon(x) > 1/2$ for large $|x|$ and $q_2(y,x) = 0$, then provided that $|x|$ is large enough
\begin{equation*}
    \alpha(x,y) \le \exp(x^2-y^2)\frac{q_1(y,x)}{0.5q_1(x,y)} = 2\exp(x^2-y^2) \frac{\exp\left(-(x-y/2)^2/2\right) }{\exp\left(-(y-x/2)^2/2\right)} = 2\exp\left(\frac{5}{8}x^2-\frac{5}{8}y^2\right)
\end{equation*}
for $x\notin (\exp(y^2),\exp(y^2)+1)$.  Noting that $V(x) = \exp(x^2/4)$ then this implies
\begin{equation*}
    \limsup_{|x|\to\infty} \frac{PV(x)}{V(x)}
    \le 
    \limsup_{|x|\to \infty}\left[(1-\varepsilon(x))\frac{Q_1V(x)}{V(x)} + 2\exp\left(\frac{5}{8}x^2-\frac{5}{8}\exp(2x^2)\right)\varepsilon(x)\frac{V(\exp(x^2)+1)}{V(x)}\right],
\end{equation*}
which upon simplification can be written
\begin{equation}
\label{eq:pv_expression}
    \limsup_{|x|\to\infty}\left[\left(1-\frac{1}{2+|x|}\right)\frac{Q_1V(x)}{V(x)} + \frac{2}{2+|x|} \exp\left(\frac{3}{8}x^2 - \frac{3}{8}\exp(2x^2)+\frac{1}{2}\exp(x^2)+\frac{1}{4}\right)\right].
\end{equation}
% We choose $V(x) =\exp(x^2/4)$. Note that
% \begin{equation*}
%     \int V(x) \pi(x) dx = c\int \exp\left(\frac{x^2}{4}\right)\exp(-x^2) dx  = c\int\exp\left(-\frac{3x^2}{4}\right)dx <\infty
% \end{equation*}
The second term on the right-hand side asymptotes to $0$ owing to the dominant $- 3\exp(2x^2)/8$ term inside the exponential.  Turning to the first we have
\begin{align*}
    \limsup_{|x|\to\infty}\frac{Q_1V(x)}{V(x)} 
    &= 
    \limsup_{|x|\to\infty}\frac{\mathbb{E}[\exp((x/2+Z)^2/4)]}{\exp(x^2/4)}
    = \limsup_{|x|\to\infty}\exp(-3x^2/16)\mathbb{E}\left[\exp\left(\frac{xZ}{4}+\frac{Z^2}{4}\right)\right].
\end{align*}
Expanding the expectation gives $\mathbb{E}\left[\exp\left(xZ/4+Z^2/4\right)\right] = c'\exp\left(x^2/16\right)$ for some $c'<\infty$, meaning
    % \int \frac{1}{\sqrt{2\pi}}\exp\left(\frac{xz}{4}+\frac{z^2}{4}-\frac{z^2}{2}\right)dz\\
    % & = \exp\left(\frac{x^2}{16}\right)\int\frac{1}{\sqrt{2\pi}}\exp\left(-\frac{(z-x/2)^2}{4}\right)dz 
    %= 
%\end{align*}
\begin{equation*}
    \lim_{|x|\to\infty}\frac{Q_1V(x)}{V(x)}  = \lim_{|x|\to\infty}c'\exp\left(-\frac{x^2}{8}\right)=0.
\end{equation*}
Substituting into \eqref{eq:pv_expression} therefore gives
\begin{equation*}
    \limsup_{|x|\to\infty} \frac{PV(x)}{V(x)}
    \le 
    \limsup_{|x|\to\infty}\frac{Q_1V(x)}{V(x)} = 0,
\end{equation*}
from which the result follows.
\end{proof}

\begin{lemma}
\label{lem:ar1_vgrowth}
If $Q_1(x,\cdot) = N(x/2, 1)$, and $V : \mathbb{R} \to [1, \infty)$ satisfies
\begin{equation}\label{eq:drift}
    Q_1 V(x) \leq \rho V(x), \qquad |x| > x_0
\end{equation}
for some $\rho \in (0,1)$ and $x_0 < \infty$, then there exist constants $c > 0$ and $\eta > 0$ such that $V(x) \geq c\, |x|^{\eta}$ for all sufficiently large $|x|$.
\end{lemma}

\begin{proof}
Let $C = [-x_0, x_0]$ and $\tau_C := \inf\{n \geq 1 : X_n \in C\}$. Set
$\varepsilon := -\log \rho > 0$.  Following the proof of Theorem~3.3 of \cite{mengersen1996rates}, it must be that any valid $V$ satisfies
\begin{equation}\label{eq:logV}
    \log V(x) \geq \varepsilon\mathbb{E}_x[\tau_C]
\end{equation}
for all sufficiently large $|x|$.\footnote{Note an inconsequential typographical error in equation (31) of \cite{mengersen1996rates}, in which the bound is reported as $\mathbb{E}_x[\tau_C]/\varepsilon$.} It therefore remains to obtain a lower bound on
$\mathbb{E}_x[\tau_C]$.  Set $n(x) := \lfloor (\log_2 x)/2 \rfloor$ for some $x > x_0$. Since
$\mathbb{P}_x(\tau_C > k)$ is non-increasing in $k$,
\[
    \mathbb{E}_x[\tau_C]
    = \sum_{k=0}^\infty \mathbb{P}_x(\tau_C > k)
    \geq n(x) \cdot \mathbb{P}_x\left(\tau_C > n(x)\right).
\]
Note that if $X_0 = x$ then $X_k \sim N(x/2^k,
\sigma_k^2)$ with $\sigma_k^2 < 4/3$. For each $k \leq n(x)$ the mean
of $X_k$ satisfies $x/2^k \geq x/2^{n(x)} \geq \sqrt{x}$, so using a Gaussian tail bound
\[
    \mathbb{P}_x(X_k \in C) \leq \Phi \left(\frac{x_0 - \sqrt{x}}{2/\sqrt{3}}\right)
    \leq e^{-c_0 x}
\]
for some constant $c_0 > 0$ and all sufficiently large $x$. A union bound over
the $n(x) = O(\log x)$ steps gives $\mathbb{P}_x(\tau_C \leq n(x)) \leq n(x)e^{-c_0 x}
\to 0$, so that $\mathbb{P}_x(\tau_C > n(x)) \geq 1/2$ for all sufficiently large $x$.
Therefore
\[
    \mathbb{E}_x[\tau_C] \geq \frac{n(x)}{2} \geq \frac{\log_2 x}{8} = \frac{\log(x)}{8\log(2)}.
\]
Substituting into \eqref{eq:logV} gives
\[
    \log V(x) \geq \log(x^\eta),
\]
with $\eta := \varepsilon/(8\log (2)) > 0$ and therefore $V(x) \geq x^{\eta}$.  An analogous argument holds when $x < -x_0$, which completes the proof.
\end{proof}


%%%%%%%%%%%%%% Section 3 %%%%%%%%%%%%%%%%%%%%%%%
\section{Proofs of Section \ref{sec: RWM and GWM on 1D R}}
\label{app:proofs_1D_RWM_GWM}


%%%%%%%%%%%%%%%%%%% 3.1 %%%%%%%%%%%%%%%%%%
\subsection{Proofs of Subsection \ref{subsec:heavy_tailed_target}}
\label{app:proofs_gwm_polynomialtails}

%\subsection{Lower bound of the polynomial convergence rate of GWM}
%\label{subsubsec: lower bound of polynomial convergence rate, GWM}
%We first give a proposition related to the lower bound of the polynomial convergence rate of GWM with finite variance proposal, together with the proof.  

\begin{lemma}
\label{lem:lowerbound_polynomialtails}
    If $\pi$ satisfies Assumption \ref{ass:heavy-tailed}, then the polynomial convergence rate of the guided walk Metropolis is at least $r$.
\end{lemma}

\begin{proof}
We prove the result by constructing a drift condition of the form
\[
PV(x) \leq V(x) - \gamma V(x)^{\alpha} + b\mathbb{I}_C(x)
\]
for some Lyapunov function $V:\mathbb{R} \to [1,\infty)$, $0 < b,\gamma < \infty$, a small set $C$ and $0 < \alpha < 1$.  Polynomial ergodicity with rate at least $\alpha/(1-\alpha)$ in total variation distance then follows from Theorem 3.6 of \citep{jarner2002polynomial}.

Consider the Lyapunov function 
\begin{equation*}
    V(x,p) = \exp(\delta \operatorname{sgn}(x)p) \pi(x)^{-\beta}
\end{equation*}
where $\delta>0$ and $1/
(r+1)<\beta<1$. Let $C = [-x_0,x_0]$ be a small set, where $x_0$ is chosen sufficiently large such that for all $|x|>x_0$, the acceptance region is $A(x) = [-|x|,|x|]$ and the rejection region is $R(x) = (-\infty,-|x|)\cup (|x|,\infty)$. Since $Z$ is Gaussian, writing $q(z)$ to denote its density we have
\begin{equation*}
    \mu = \int_0^\infty z q(z) dz < \infty, \quad \sigma^2 = \int_0^\infty z^2 q(z) dz < \infty.
\end{equation*}
Note that the Lyapunov function $V(x,p)$ shows two different patterns in the exponential term, and we will discuss these two cases separately.

\textit{Case 1:} $\operatorname{sgn}(x)p > 0$.  Here $V(x,p) = e^\delta \pi(x)^{-\beta} = e^\delta |x|^{(r+1)\beta}$. The transition operator gives
\begin{equation}
    PV(x,p) = 2\int_0^\infty V(x+pz,p)\frac{\pi(x+pz)}{\pi(x)}q(z)dz + 2\int_0^\infty V(x,-p)\left(1-\frac{\pi(x+pz)}{\pi(x)} \right)q(z)dz
\end{equation}
Since $\operatorname{sgn}(x)p>0$, we have $\operatorname{sgn}(x+pz)p>0$ for all $z>0$. For the rejection probability, by Taylor's theorem with explicit remainder there exists $\xi \in (0, z/|x|)$ such that
\begin{equation*}
    1 - \left(1+\frac{z}{|x|}\right)^{-(1+r)} = \frac{(1+r)z}{|x|} - \frac{(1+r)(2+r)}{2}\frac{z^2}{|x|^2}\left(1+\xi\right)^{-(3+r)}.
\end{equation*}
Since $(1+\xi)^{-(3+r)} \leq 1$ for $\xi > 0$, the remainder is bounded in absolute value by $\frac{(1+r)(2+r)}{2}\frac{z^2}{|x|^2}$. Integrating term by term against $q(z)$ is therefore justified by the finite second moment $\sigma^2 < \infty$, giving
\begin{equation*}
    \int_0^\infty \left(1-\frac{\pi(x+pz)}{\pi(x)}\right)q(z)dz = \frac{(1+r)\mu}{|x|} - R(x),
\end{equation*}
where the remainder satisfies
\begin{equation*}
    0 \leq R(x) \leq \frac{(1+r)(2+r)\sigma^2}{2|x|^2}.
\end{equation*}
Thus,
\begin{equation}
\begin{aligned}
    \frac{PV(x,p)-V(x,p)}{V(x,p)} 
    &= 2\int_0^\infty \left[\left(\frac{\pi(x+pz)}{\pi(x)}\right)^{1-\beta}-1\right]q(z)dz + 2e^{-2\delta}\int_0^\infty \left(1-\frac{\pi(x+pz)}{\pi(x)}\right) q(z) dz\\
    &\le 2\int_0^\infty\left[\left(1+\frac{z}{|x|}\right)^{(r+1)(\beta-1)}-1\right]q(z)dz + 2e^{-2\delta}\left[\frac{(1+r)\mu}{|x|} - R(x)\right].
\end{aligned}
\end{equation}
Applying Taylor's theorem to the first integrand as before gives
\begin{align}
    \frac{PV(x,p)-V(x,p)}{V(x,p)}  &\le -\frac{2(r+1)(1-\beta)\mu}{|x|} + \frac{c'\sigma^2}{|x|^2} + \frac{2e^{-2\delta}(1+r)\mu}{|x|}\\
    &= \frac{2(r+1)\mu(\beta - 1 + e^{-2\delta})}{|x|} + \frac{c'\sigma^2}{|x|^2}.
\end{align}
Now choose $\delta$ large enough that $e^{-2\delta} < 1-\beta$, which is possible since $\beta < 1$, so that the coefficient $\beta - 1 + e^{-2\delta} < 0$. Setting $c_1 = 2(r+1)\mu(1-\beta-e^{-2\delta}) > 0$, and choosing $|x|$ large enough that $c'\sigma^2/|x| < c_1/2$, gives
\begin{align}
    \frac{PV(x,p)-V(x,p)}{V(x,p)} \leq -\frac{c}{|x|}
\end{align}
for some $c > 0$, as required.

\textit{Case 2:} $\operatorname{sgn}(x)p<0$.  Now $V(x,p) =e^{-\delta} \pi(x)^{-\beta} = e^{-\delta} |x|^{(r+1)\beta}$. The transition operator gives
\begin{equation}
\begin{aligned}
    PV(x,p) 
    &= 2\int_0^{|x|} V(x+pz,p) q(z) dz + 2\int_{|x|}^{2|x|} V(x+pz,p) q(z)dz\\
    &\quad + 2\int_{2|x|}^\infty V(x+pz,p)\frac{\pi(x+pz)}{\pi(x)}q(z)dz \\
    &\quad + 2\int_{2|x|}^\infty V(x,-p) \left(1-\frac{\pi(x+pz)}{\pi(x)}\right)q(z)dz.
\end{aligned}
\end{equation}
For $z\in(0,|x|)$, we have $\operatorname{sgn}(x+pz)p <0$, so $V(x+pz,p) = e^{-\delta} |x+pz|^{(r+1)\beta}$. For $z \in (|x|,2|x|)$, we have $\operatorname{sgn}(x+pz)p>0$, so $V(x+pz,p) = e^\delta |x+pz|^{(r+1)\beta}$. Therefore
\begin{equation}
\begin{aligned}
    \frac{PV(x,p) - V(x,p)}{V(x,p)}
    &= 2\int_0^{|x|} \left[\left|\frac{x+pz}{x}\right|^{(r+1)\beta}-1\right]q(z) dz + 2e^{2\delta} \int_{|x|}^{2|x|} \left|\frac{x+pz}{x}\right|^{(r+1)\beta}q(z)dz\\
    &\quad + 2e^{2\delta}\int_{2|x|}^\infty \left|\frac{x+pz}{x}\right|^{(r+1)\beta}\frac{\pi(x+pz)}{\pi(x)}q(z)dz\\
    &\quad + 2e^{2\delta}\int_{2|x|}^\infty \left(1- \frac{\pi(x+pz)}{\pi(x)}\right)q(z) dz -2\int_{|x|}^\infty q(z)dz\\
    &\le 2\int_0^{|x|}\left[\left|\frac{x+pz}{x}\right|^{(r+1)\beta}-1\right]q(z) dz + 2(e^{2\delta}-1) \mathbb{P}(Z > |x|).
\end{aligned}
\end{equation}
Applying Taylor's theorem with explicit remainder to the first integrand, for $z \in (0,|x|)$ and $\operatorname{sgn}(x)p < 0$ there exists $\xi \in (0, z/|x|)$ such that
\begin{equation*}
    \left(1-\frac{z}{|x|}\right)^{(r+1)\beta} - 1 = -(r+1)\beta\frac{z}{|x|} + \frac{(r+1)\beta((r+1)\beta-1)}{2}\frac{z^2}{|x|^2}(1-\xi)^{(r+1)\beta - 2}.
\end{equation*}
Since $(1-\xi)^{(r+1)\beta-2} \leq 1$ for $\xi \in (0,1)$ and $(r+1)\beta > 1$ by assumption, the remainder is bounded in absolute value by $c'' z^2/|x|^2$ where $c'' = (r+1)\beta((r+1)\beta-1)/2 > 0$. Integrating term by term against $q(z)$ over $(0,|x|)$ is therefore justified by $\sigma^2 < \infty$, giving
\begin{align}
    \frac{PV(x,p) - V(x,p)}{V(x,p)} &\le -\frac{2(r+1)\beta}{|x|}\int_0^{|x|}zq(z)dz + \frac{c''\sigma^2}{|x|^2} + 2(e^{2\delta}-1)\mathbb{P}(Z>|x|).
\end{align}
Since $q(z)$ is Gaussian, $\mathbb{P}(Z > |x|)$ decays super-exponentially in $|x|$, so for any fixed $\delta$ we have $2(e^{2\delta}-1)\mathbb{P}(Z>|x|) = o(1/|x|)$ as $|x| \to \infty$. Moreover $\int_0^{|x|}zq(z)dz \to \mu$ as $|x|\to\infty$, so for sufficiently large $|x|$ the dominant term is $-2(r+1)\beta\mu/|x|$, and the same conclusion as in case 1 holds, namely
\begin{align}
    \frac{PV(x,p) - V(x,p)}{V(x,p)} \leq -\frac{c}{|x|}
\end{align}
for some $c > 0$.

Combining both cases shows that outside $C$ there exists $c^*>0$ such that the polynomial drift condition
\begin{equation}
    PV(x,p) - V(x,p) \le -c^*V(x,p)^{((r+1)\beta-1)/[(r+1)\beta]}
\end{equation}
holds. The polynomial rate is therefore at least
\begin{equation}
    \frac{((r+1)\beta-1)/[(r+1)\beta]}{1-((r+1)\beta-1)/[(r+1)\beta]} = (r+1)\beta -1.
\end{equation}
For any $\varepsilon > 0$, choosing $\beta = 1-\varepsilon \in (1/(r+1), 1)$ gives a valid Lyapunov function and a polynomial rate of at least $(r+1)(1-\varepsilon) - 1 = r - (r+1)\varepsilon$. Since this holds for all $\varepsilon > 0$, we conclude that the polynomial rate is at least $\sup_{\varepsilon > 0}\{r - (r+1)\varepsilon\} = r$, which completes the proof.
\end{proof}


%\subsection{Upper bound of the polynomial convergence rate of GWM}
%\label{subsubsec: upper bound of polynomial convergence rate, GWM}
%Correspondingly, in this section we give and prove the proposition related to the upper bound of the polynomial convergence rate of GWM with finite variance proposal.

\begin{lemma}
\label{lem:upperbound_polynomialtails}
    If $\pi$ satisfies Assumption \ref{ass:heavy-tailed}, then the polynomial convergence rate of the guided walk Metropolis with finite variance proposal is at most $r$.
\end{lemma}

\begin{proof}
Suppose that the algorithm begins from $(x,p) = (0,1)$. For any $n>0$, we can choose $k>0$ such that $kn>K$. Define the set $A_n = \{x:x>kn\}$, and let $\mu$ be the joint invariant measure of $X,P$. For notational simplicity we write $P^n((x,p),A)$ for $P^n((x,p),(A,\{+1,-1\}))$, and note that $\mu((A,\{+1,-1\})) = \pi(A)$.  By the definition of total variation distance, we have
\begin{align}
    \|P^n((0,1),\cdot) - \mu\|_{TV} &\ge |P^n((0,1),A_n) - \pi(A_n)|.
\end{align}    
To bound $P^n((0,1),A_n)$, we observe that if the chain is never rejected in $n$ steps, it continuously moves in the positive direction. By a stochastic comparison argument the probability that the chain reaches $A_n$ is at most the probability of reaching $A_n$ without any rejections. Therefore,
\begin{equation*}
    P^n((0,1),A_n) \le \mathbb{P}\left(\sum_{i=1}^n |z_i|>kn\right), \quad \text{where } z_i \stackrel{\text{i.i.d.}}{\sim} N(0,\sigma^2) \text{for} \, i=1, \cdots ,n.
\end{equation*}
Since each $|z_i|$ follows a half-normal distribution with tail decay at least as fast as a Gaussian, $|z_i|$ is sub-Gaussian. Let $\mu = \mathbb{E}[|z_i|] = \sigma\sqrt{2/\pi}$ and $\sigma_H^2 = \text{Var}(|z_i|)$. Then for all $\lambda>0$,
\begin{equation*}
    \mathbb{E}\left[e^{\lambda(\sum_{i=1}^n |z_i| - n\mu)}\right]\le \exp\left(\frac{n\lambda^2\sigma_H^2}{2}\right).
\end{equation*}
Applying Markov's inequality to the exponential function, we obtain for all $\lambda>0$
\begin{equation}
\begin{aligned}
    \mathbb{P}\left(\sum_{i=1}^n|z_i| >kn\right) &= \mathbb{P}\left(\sum_{i=1}^n|z_i| - n\mu>(k-\mu)n\right)\\
    &= \mathbb{P}\left(\exp\left(\lambda \sum_{i=1}^n(|z_i| - \mu)\right) >\exp(\lambda(k-\mu)n)\right)\\
    &\le \frac{\mathbb{E}\left[\exp\left(\lambda \sum_{i=1}^n(|z_i| - \mu)\right)\right]}{\exp(\lambda(k-\mu)n)}\\
    &\le \exp\left(\frac{n\lambda^2\sigma_H^2}{2} - \lambda n (k-\mu)\right).
\end{aligned}
\end{equation}
Since this holds for all $\lambda>0$, we can minimize over $\lambda$ by choosing $\lambda^* = (k-\mu)/\sigma_H^2$, giving
\begin{align}
    P^n((0,1),A_n) \le  \mathbb{P}\left(\sum_{i=1}^n|z_i| >kn\right)\le \exp\left(-\frac{n(k-\mu)^2}{2\sigma_H^2}\right).
\end{align}
Since $kn>K$ by construction, for $x>kn$ we have $\pi(x) = C_0 x^{-(1+r)}$. Combining gives
\begin{equation}
    \pi(A_n) = \int_{kn}^\infty \frac{C_0}{x^{1+r}}  dx = \frac{C_0}{r(kn)^r} = \frac{C'}{n^r}, \quad \text{where } C' = C_0/(rk^r),
\end{equation}  
which implies that for all $n\ge n_0$, for suitably large $n_0$ and constant $k>0$
\begin{equation}
    |\pi(A_n) - P^n((0,1),A_n)| \ge \left|\frac{C'}{n^r} - \exp\left(-\frac{n(k-\mu)^2}{2\sigma_H^2}\right)\right| \ge \frac{c}{n^r}
\end{equation}
where $c>0$ is a constant. The last inequality holds because for any fixed $r>0$, we can choose large $k=k(r)$, and $C' = C_0/(rk^r):= C'(r)$ such that both relate to $r$ and $\exp(-n(k-\mu)^2/(2\sigma_H^2)) < C'/(2n^r)$ for all $n\ge n_0$ for suitably large $n_0$, which implies the bound with $c = C'/2$.  We have therefore established that
\begin{equation}
    \|P^n((0,1),\cdot) - \pi(\cdot)\|_{TV} \ge \frac{c}{n^r}, \quad n\ge n_0,
\end{equation}
from which the result follows.
\end{proof}


%%%%%%%%%%%%%%%%%% 3.2 %%%%%%%%%%%%%%%%%%%%%%%%%
\subsection{Proofs of Subsection \ref{subsec:light_tailed_target}}
\label{app:proofs_gwm_lighttails}

\begin{lemma}
\label{lem:convergence_alpha_light_tails}
If $\pi$ satisfies Assumption \ref{ass:strict_convexity}, then the (Gaussian) random walk Metropolis satisfies
\begin{equation*}
    \lim_{x\to \infty}\int_{x}^{\infty} \alpha(x,y)Q(x,dy) = \lim_{x\to -\infty}\int_{-\infty}^{x} \alpha(x,y)Q(x,dy) = 0.
\end{equation*}
\end{lemma}

\begin{proof}
Considering $[x, \infty)$ for $x>0$
\begin{equation*}
\begin{aligned}
    \int_x^\infty \alpha(x,y)Q(x,dy) = \int_{0}^{\infty} \alpha(x, x+\epsilon z)\phi(z) dz
    = 
    \int_0^{\infty} \min(1, \exp(U(x) - U(x+\epsilon z)))\phi(z)dz
\end{aligned}
\end{equation*}
where $\phi(z)$ is the standard Normal density. The mean value theorem states that
\begin{equation*}
    U(x+\epsilon z) - U(x) \geq \epsilon z \inf_{t \in [x, x+\epsilon z]} U'(t) = \epsilon z U'(x)
\end{equation*}
where the last equality follows from monotonicity of $U'$.  By assumption $U(x+\epsilon z) -  U(x) \to \infty$ as $x \to \infty$ because $U'(t) \to \infty$, giving the point-wise limit
\begin{equation*}
    \lim_{x\to \infty}\alpha(x, x+\epsilon z) = 0
\end{equation*}
for any fixed $z > 0$. Since $\alpha(x,x+\epsilon z) \leq 1$ for any $z$, then the bounded convergence theorem can be used to conclude.  For the tail region $(-\infty, x]$ when $x \to -\infty$ an identical argument can be applied, which completes the proof.
\end{proof}

% Now we assume that we have a sequence $\{x_n\}$ that diverges to $\infty$ and we denote $f_n(z) = \alpha(x_n, x_n + \epsilon z)$, and $f(z) = 0$ as a constant function. Easily we have $f_n \to f$ when $n \to \infty$, moreover for another constant function $g(z)=1$ we have $|f_n| \leq g$, and $\int_0^{\infty}g(z)\varphi(z)dz = 1/2 < \infty$. Then we can apply the dominated convergence theorem to have
% \begin{equation*}
%     \lim_{n\to \infty}\int_{0}^{\infty} f_n(z) \varphi(z)dz = \int_{0}^{\infty} \lim_{n \to \infty}f_n(z) \varphi(z)dz = \int_{0}^\infty f(z)\varphi(z)dz = 0
% \end{equation*}
% which is equivalent to $\lim_{x\to \infty}I_1(x) = \int_0^{\infty} \lim_{x\to \infty}\alpha(x, x+\epsilon z)\varphi(z) = 0$. 

% For the tail region $(-\infty, x]$ for $x\leq 0$, we apply the mean value theorem in the way that, for $u>0$
% \begin{equation*}
%     U(x) - U(x-\epsilon u) \leq \epsilon u \sup_{t \in [x-\epsilon u, x]}U'(t) = \epsilon u U'(x)
% \end{equation*}
% so with $U'(x) \to -\infty$ when $x\to -\infty$, we clearly have $U(x) - U(x-\epsilon u) \to -\infty \implies \alpha(x, x-\epsilon u) \to 0$ when $x \to -\infty$. Therefore,
% \begin{equation*}
% \begin{aligned}
%     I_2(x) = \int_{-\infty}^{x}\alpha(x,y)q(x,dy) &= \int_{-\infty}^{0} \alpha(x, x+\epsilon z)\phi(z) dz\\
%     &= \int_0^{\infty} \alpha(x, x-\epsilon z)\phi(z)dz \to 0 \text{ when $x \to -\infty$}
% \end{aligned}
% \end{equation*}
% again by using the dominated convergence theorem.
% \end{proof}

% \begin{proof}[Proof of Lemma \ref{lemma: stability of momentum in the tail region}]
% Before the proof, note that for the general light-tailed target, we have $U'(x)$ being strictly monotonic increasing, then if we have $K > 0$ such that $U'(K) > 0$ and $U'(-K) \leq 0$, automatically we shall have $U(x)$ being strictly monotonic increasing for $[K, \infty)$, and monotonic decreasing for $(-\infty, -K]$.

% We now prove by induction on $t$. For the base case $t=1$, WLOG assume we start with $x>0$ and $x\to \infty$, then $P_0 = -\mathrm{sgn}(x) = -1$. We will carry the LGWM proposal when $P_0 z \geq 0 \iff z \leq 0$, where $x'=w'= x + \epsilon z \leq x$. Since $x \to \infty$, for any $z$, $x + \epsilon z$ eventually lies in the region where $U$ is strictly increasing. i.e. with $K$ such that $U$ is monotonic increasing in $[K, \infty)$, and $x \geq K$, the probability for $w'=x+\epsilon z \in [K, \infty)$ is 
% \begin{equation*}
%     P(x + \epsilon z \geq K) = P\left(z \geq -\frac{K-x}{\epsilon}\right) \to 1
% \end{equation*}
% Thus, under $P_0 z \geq 0$ we have $x'=w'>0$ with probability tending to 1 when $x\to \infty$, and we will have $P_1 = P_0 = -1 = -\mathrm{sgn}(W_1)$ if we accept the move $W_1 = w'>0$, with acceptance probability
% \begin{equation*}
%     \alpha(x, w') = \min(1, \exp(U(x) - U(w'))) = 1
% \end{equation*}
% because $w'\leq x \implies U(w')\leq U(x)$ with probability tending to 1 when $x\to \infty$. On the other hand we shall keep $W_1 = x > 0$ and $P_1 = P_0 = -1=-\mathrm{sgn}(W_1)$ when $P_0 z \leq 0$. Therefore, altogether we will have
% \begin{equation*}
%     P(A_1 \mid R) = P(P_1 = -\mathrm{sgn}(W_1) \mid R) \to 1 \quad \text{when $x \to \infty$} 
% \end{equation*}
% Now we take the inductive hypothesis that $P(A_{t} \mid E_{t-1}^c \cap \cdots \cap E_1^c \cap R) \to 1$, meaning that $P_{t} = -\mathrm{sgn}(W_{t})$ with probability tending to 1. Note that under $x\to \infty$ we shall have all intermediate states being positive with probability tending to 1, i.e. $X_{t}, W_t > 0$ for all $t \geq 0$ with very high probability. To explain this, note that the maximum total displacement from initial $x$ via $t$ steps is $\sum_{i=1}^t\epsilon|z_i|$, meaning that the smallest value after $t$ steps is $x - \sum_{i=1}^t \epsilon|z_i|$. With finite $t$ and i.i.d. $z_i \sim N(0,1)$, when $x\to \infty$ by Markov inequality we have
% \begin{equation*}
%     P\left(x - \sum_{i=1}^t \epsilon|z_i| \leq 0\right) =  P\left( \sum_{i=1}^t \epsilon|z_i| \geq x\right) \leq \frac{\epsilon \sum_{i=1}^t \mathbb{E}(|z_i|)}{x} \to 0
% \end{equation*}
% showing that even for the smallest possible value, when $x\to \infty$ we have that smallest value being positive with very low probability, hence $P_{t} = -1$ with probability tending to 1. 

% We examine the transition to step $t+1$ given the state at $t$. When under $X_{t} = W_{t}$ within $R$, the coupling says that if $P_{t}z > 0$ then $x'= w' = W_{t} + \epsilon z \leq W_{t}$, and within $R$ we have $W_{t} = X_{t} \geq K$. With the same logic as the base case we know that the probability for $w' \in [K, \infty)$ lies within the monotonicity interval approaches 1 when $K$ is large enough, so with $U(w') \leq U(X_{t})$ we will have
% \begin{equation*}
%     \alpha(X_{t}, w') = \min(1, \exp(U(X_t) - U(w'))) = 1
% \end{equation*}
% meaning that under $P_t z > 0$ we will accept the move ($W_{t+1} =w'$) and have $P_{t+1} = P_t = -1$ with probability 1. Also with $X_t=W_t > 0$ we shall have $P(w'>0) \to 1$ when the initial state $x\to \infty$, therefore $P_{t+1} = -\mathrm{sgn}(W_{t-1})$ under $P_t z > 0$. On the other hand, for $P_t z \leq 0$ we directly keep $W_{t+1} = W_t > 0$ and $P_{t+1} = P_t = -1 = -\mathrm{sgn}(W_{t+1})$. Altogether the claim is true for $t+1$ as well, i.e.
% \begin{equation*}
%     P\left(A_{t+1} \mid E_{t}^c \cap E_{t-1}^c \cap \cdots \cap E_1^c \cap R\right) \to 1 
% \end{equation*}
% Hence we prove Lemma \ref{lemma: stability of momentum in the tail region} by induction.
% \end{proof}

% \begin{proof}[Proof of Lemma \ref{lemma:out_of_R_zero_prob}]
% We shall prove this proposition by induction on $t$. When $t=1$, with $X_1 = X_0 + \epsilon z = x+\epsilon z, z \sim N(0,1)$ we write
% \begin{equation*}
% \begin{aligned}
%     P(|X_1| \leq K \mid X_0 = x) &= P(-K \leq x+\epsilon z \leq K)\\
%     &= \Phi\left(\frac{K-x}{\epsilon}\right) - \Phi\left(\frac{-K-x}{\epsilon}\right)
% \end{aligned}
% \end{equation*}
% Now when $x \to \pm \infty$ we obviously have $P(|X_1| \leq K \mid X_0 = x) \to 0$, with fixed $K>0$ and stepsize $\epsilon > 0$. Now we assume the claim holds at $t$, and define
% \begin{equation}
%     f(u) = P(|X_1|\leq K \mid X_0 = u)
% \end{equation}
% From the base case we know $f(u) \to 0$ when $u \to \pm \infty$, which equivalently means
% \begin{equation}
% \label{eqn: appx proof limit form}
%     \forall \kappa > 0, \exists \delta = \delta(\kappa) > 0, \forall u \text{ s.t. } |u| > \delta, f(u) < \kappa
% \end{equation}
% and now we fix that $\kappa$ (so as related $\delta = \delta(\kappa)$). At $t+1$, we can use $f$ to write (TIME-HOMOGENEOUS ASSUMPTION...)
% \begin{equation}
% \begin{aligned}
% &P(|X_{t+1}| \leq K\mid X_0=x) = \mathbb{E}_{X_t}(f(X_t)\mid X_0 = x)\\
% &= \mathbb{E}(f(X_t)\mathbf{1}(|X_t| > \delta(\kappa))\mid X_0=x) + \mathbb{E}(f(X_t)\mathbf{1}(|X_t| \leq \delta(\kappa))\mid X_0=x)
% \end{aligned}
% \end{equation}
% Note that by \eqref{eqn: appx proof limit form} we see $\mathbb{E}(f(X_t)\mathbf{1}(|X_t| > \delta(\kappa))\mid X_0=x) < \kappa$ for any arbitrary $\kappa>0$ (as here we have $|X_t| > \delta(\kappa) \implies f(X_t) < \kappa$), while
% \begin{equation*}
%     \mathbb{E}(f(X_t)\mathbf{1}(|X_t| \leq \delta(\kappa))\mid X_0=x) \leq \left[\sup_{|u| \leq \delta(\kappa)} f(u)\right] P(|X_t|\leq \delta(\kappa) \mid X_0=x) \to 0
% \end{equation*}
% based on the inductive hypothesis that $P(|X_t|\leq K \mid X_0=x) \to 0$ for any arbitrary $K > 0$, and here we just set $K=\delta(\kappa)$. Together, we get $P(|X_{t+1}| \leq K\mid X_0=x) \to 0$ when $x\to \pm \infty$, so our claim also holds at $t+1$, and then by induction, the claim is proved.
% \end{proof}



% % We shall prove this proposition by induction on $t$. When $t=1$, with $X_1 = X_0 + \epsilon z = x+\epsilon z, z \sim N(0,1)$ we write
% % \begin{equation*}
% % \begin{aligned}
% %     P(|U'(X_1)| \leq K \mid X_0 = x) &= P(-K \leq U'(x+\epsilon z) \leq K)\\
% %     &= P(k_1 \leq x+\epsilon z \leq k_2)\\
% %     &= \Phi\left(\frac{k_2-x}{\epsilon}\right) - \Phi\left(\frac{k_1-x}{\epsilon}\right)
% % \end{aligned}
% % \end{equation*}
% % where $k_1, k_2$ are defined to satisfy
% % \begin{equation*}
% %     U'(k_1) = -K, \quad U'(k_2) = K, \quad k_1 \leq k_2
% % \end{equation*}
% % We know such $k_1, k_2$ exist as for general light-tailed target $U$, we have $U'(x)$ is strictly monotonic increasing and continuous (POSSIBLE GRAPHIC ILLUSTRATION?). Now when $x \to \pm \infty$ we obviously have $P(|U'(X_1)| \leq K \mid X_0 = x) \to 0$, with fixed $K>0$ and stepsize $\epsilon > 0$. 

% % Now we assume the claim holds at $t$, and define
% % \begin{equation}
% %     f(u) = P(|U'(X_1)|\leq K \mid X_0 = U'(u))
% % \end{equation}
% % From the base case we know $f(u) \to 0$ when $u \to \pm \infty$ (as $|U'(x)|\to \infty$ when $|x|\to \infty$), which equivalently means
% % \begin{equation}
% % \label{eqn: appx proof limit form}
% %     \forall \kappa > 0, \exists \delta = \delta(\kappa) > 0, \forall u \text{ s.t. }|U'(u)| > \delta, f(u) < \kappa
% % \end{equation}
% % and now we fix that $\kappa$ (so as related $\delta = \delta(\kappa)$). At $t+1$, we can use $f$ to write (TIME-HOMOGENEOUS ASSUMPTION...)
% % \begin{equation}
% % \begin{aligned}
% % &P(|U'(X_{t+1})| \leq K\mid X_0=x) = \mathbb{E}_{X_t}(f(X_t)\mid X_0 = x)\\
% % &= \mathbb{E}(f(X_t)\mathbf{1}(|U'(X_t)| > \delta(\kappa))\mid X_0=x) + \mathbb{E}(f(X_t)\mathbf{1}(|U'(X_t)| \leq \delta(\kappa))\mid X_0=x)
% % \end{aligned}
% % \end{equation}
% % Note that by \eqref{eqn: appx proof limit form} we see $\mathbb{E}(f(X_t)\mathbf{1}(|U'(X_t)| > \delta(\kappa))\mid X_0=x) < \kappa$ for any arbitrary $\kappa>0$ (as here we have $|U'(X_t)| > \delta(\kappa) \implies f(X_t) < \kappa$), while
% % \begin{equation*}
% %     \mathbb{E}(f(X_t)\mathbf{1}(|X_t| \leq \delta(\kappa))\mid X_0=x) \leq \left[\sup_{|u| \leq \delta(\kappa)} f(u)\right] P(|U'(X_t)|\leq \delta(\kappa) \mid X_0=x) \to 0
% % \end{equation*}
% % based on the inductive hypothesis that $P(|U'(X_t)|\leq K \mid X_0=x) \to 0$ for any arbitrary $K > 0$, and here we just set $K=\delta(\kappa)$. Together, we get $P(|X_{t+1}| \leq K\mid X_0=x) \to 0$ when $x\to \pm \infty$, so our claim also holds at $t+1$, and then by induction, the claim is proved.
% % \end{proof}



% % \section{Supplementary material}

% % \subsection{Additional lemmas, propositions, remarks and proofs}
% % \label{appx subsec: additional proofs}

% \begin{lemma}[convergence of acceptance probability at the tail, 1D canonical Gaussian target]
% \label{lemma: convergence of acceptance probability at the tail, 1D canonical Gaussian}
%     With the target density being as the 1D canonical Gaussian and coupling described in section \ref{subsubsec: first-step RWM and LGWM}, for $z>0$, the acceptance probability for RWM transition converges to 0 when $x \to \infty$.
% \end{lemma}

% \begin{proof}
% Such acceptance probability is directly computed as
% \begin{equation}
% \begin{aligned}
%      P(z>0, X \text{ transition}) &= P(X \text{ transition} \mid z>0)P(z>0)\\
%      &= \frac{1}{2}\int_x^{\infty}\alpha(x,y)q(x,dy) = \frac{1}{2}\lim_{c\to \infty}\textcircled{1}
% \end{aligned}
% \end{equation}
% where $\textcircled{1} = \int_{x}^c \alpha(x,y)q(x,dy)$ simply represents the acceptance probability at the tail region $[x, c]$. When the target density is 1D canonical Gaussian, we compute the explicit form of \textcircled{1} in terms of

% \begin{equation}
% \label{eqn: expression of textcircled 1}
% \begin{aligned}
%     &\textcircled{1} = \int_x^{c} \alpha(x, y)q(x,dy) = \int_x^c \exp \left(\frac{1}{2}(x^2 - y^2) \right) \frac{1}{\epsilon\sqrt{2\pi}}\exp\left(-\frac{1}{2\epsilon^2}(y-x)^2\right)dy\\
%     &= \frac{1}{\epsilon \sqrt{2\pi}}\int_x^c \exp\left(\frac{1}{2}x^2(1 - \epsilon^{-2})\right) \exp\left(-\frac{1}{2\frac{\epsilon^2}{1 + \epsilon^2}}\left(y - \frac{1}{\epsilon^2(1 + \epsilon^{-2})}x \right)^2\right)\exp\left(\frac{1}{2\epsilon^2(1+\epsilon^2)}x^2\right) dy\\
%     &= \frac{1}{\epsilon \sqrt{2\pi}} \exp\left(\frac{\epsilon^2}{2(1+\epsilon^2)}x^2\right) \int_{x - \frac{1}{\epsilon^2(1 + \epsilon^{-2})}x}^{c - \frac{1}{\epsilon^2(1 + \epsilon^{-2})}x}  \exp\left(-\frac{u^2}{2\frac{\epsilon^2}{1 + \epsilon^2}}\right)du \quad \text{by $u = y - \frac{1}{\epsilon^2(1 + \epsilon^{-2})}x$}\\
%     &= \frac{1}{\epsilon \sqrt{2\pi}} \exp\left(\frac{\epsilon^2}{2(1+\epsilon^2)}x^2\right) \int_{g_2(x)}^{g_1(c, x)}  \exp\left(-\frac{t^2}{2}\right)\frac{\epsilon}{\sqrt{1+\epsilon^2}}dt \quad \text{by $t = \frac{u}{\epsilon / \sqrt{1 + \epsilon^2}}$}\\
%     &= \frac{1}{\sqrt{1 + \epsilon^2}} \exp\left(\frac{\epsilon^2}{2(1+\epsilon^2)}x^2\right) \int_{g_2(x)}^{g_1(c, x)} \frac{1}{\sqrt{2\pi}}\exp\left(-\frac{t^2}{2}\right)dt\\
%     &= \frac{1}{\sqrt{1 + \epsilon^2}} \exp\left(\frac{\epsilon^2}{2(1+\epsilon^2)}x^2\right) \left[\Phi(g_1(c, x)) - \Phi(g_2(x))\right]
% \end{aligned}
% \end{equation}
% where $\Phi$ is the CDF for canonical normal distribution $N(0,1)$, and
% \begin{align}  
% \label{eqn: expression of g1}
%     g_1(c, x) &= \frac{\sqrt{1+\epsilon^2}}{\epsilon}\left(c - \frac{1}{\epsilon^2(1 + \epsilon^{-2})}x\right) = \frac{\sqrt{1+\epsilon^2}}{\epsilon}\left(c-\frac{x}{1+\epsilon^2}\right)\\
% \label{eqn: expression of g2}
%     g_2(x) &= \frac{\sqrt{1+\epsilon^2}}{\epsilon}\left(x - \frac{1}{\epsilon^2(1 + \epsilon^{-2})}x\right) = \frac{\epsilon}{\sqrt{1+\epsilon^2}}x
% \end{align}
% When $c \to \infty$ (so the tail region is $[x, \infty)$), we intuitively denote
% \begin{equation}
%     \lim_{c\to \infty}\textcircled{1} = \frac{1}{\sqrt{1 + \epsilon^2}} \exp\left(\frac{\epsilon^2}{2(1+\epsilon^2)}x^2\right) \left[1 - \Phi(g_2(x))\right]
% \end{equation}
% For \textcircled{1} evaluated in \eqref{eqn: expression of textcircled 1}, when $x \to \infty$ we will have $\lim_{x\to \infty}\textcircled{1} = 0$, which can be deduced by using the standard Gaussian tail asymptotic (Mills ratio) that $1-\Phi(t)\sim \varphi(t)/t$ when $t \to \infty$, so the exponential terms will cancel and remain a decay term $O(x^{-1})$. Therefore, the lemma is proved with
% \begin{equation*}
%     \lim_{x \to \infty} P(z>0, X \text{ transition}) = 0
% \end{equation*}
% \end{proof}


% % \begin{lemma}[convergence of acceptance probability at the tail, general light-tailed target]
% % \label{lemma: convergence of acceptance probability at the tail, general light-tailed target}
% % For the general form of light-tailed target distribution in the form of $\pi(x) \propto \exp(-U(x))$, with $U(\cdot)$ being as (i). strictly convex and (ii). diverges super-linearly. Then with the RWM proposal $q$, the acceptance probability for RWM transition at the tail region $[x, \infty)$ converges to 0 when $x \to \infty$. i.e. we shall have
% % \begin{equation*}
% %     I(x) = \int_{x}^{\infty} \alpha(x,y)q(x,dy) \to 0
% % \end{equation*}
% % \end{lemma}

% % \begin{proof}[Proof of lemma \ref{lemma: convergence of acceptance probability at the tail, general light-tailed target}]
% % We say such $U(\cdot)$ satisfies strict convexity and grows super-linearly if for all $x \in \mathcal{X} = \mathbb{R}$
% % \begin{equation*}
% %     U'(x) > 0, \quad \lim_{x \to \infty}\frac{U(x)}{x} = \infty
% % \end{equation*}
% % We claim $U'(x) \to \infty$ when $x \to \infty$, and it can be shown by contradiction: with convexity, assume $U'(x) \in (0, M), \forall x \in \mathbb{R}$ is bounded for some positive, finite $M$. Then for $0<p<x$ we have
% % \begin{equation*}
% % \begin{aligned}
% %     U(x) - U(p) = \int_p^x U'(t)dt \leq \int_p^x M dt = M(x-p) &\implies U(x) \leq Mx + U(p)\\
% %     &\implies \lim_{x\to \infty}\frac{U(x)}{x} \leq M
% % \end{aligned}
% % \end{equation*}
% % which contradicts with our setting of super-linear growth, hence by contradiction we show $U'(x) \to \infty$ when $x \to \infty$. Now with the RWM proposal $q(x, y) = N(y; x, \epsilon^2)$, we have
% % \begin{equation}
% % \begin{aligned}
% %     I(x) = \int_x^{\infty}\alpha(x,y)q(x,dy) &= \int_{0}^{\infty} \alpha(x, x+\epsilon z)\varphi(z) dz\\
% %     &= \int_0^{\infty} \min(1, \exp(U(x) - U(x+\epsilon z)))\varphi(z)dz
% % \end{aligned}
% % \end{equation}
% % where $\varphi(z)$ is the density for $N(0,1)$. The mean value theorem states that
% % \begin{equation*}
% %     U(x+\epsilon z) - U(x) \geq \epsilon z \inf_{t \in [x, x+\epsilon z]} U'(t)
% % \end{equation*}
% % hence for fixed stepsize $\epsilon$, when $x \to \infty$ we have $U(x+\epsilon z) -  U(x) \to \infty$ due to $U'(t) \to \infty$. This indicates our acceptance probability at the tail satisfies
% % \begin{equation*}
% %     \lim_{x\to \infty}\alpha(x, x+\epsilon z) = \min(1, \lim_{x\to \infty} \exp(U(x) - U(x+\epsilon z))) = 0
% % \end{equation*}
% % Now we assume that we have a sequence $\{x_n\}$ that diverges to $\infty$ and we denote $f_n(z) = \alpha(x_n, x_n + \epsilon z)$, and $f(z) = 0$ as a constant function. Easily we have $f_n \to f$ when $n \to \infty$, moreover for another constant function $g(z)=1$ we have $|f_n| \leq g$, and $\int_0^{\infty}g(z)\varphi(z)dz = 1/2 < \infty$. Then we can apply the dominated convergence theorem to have
% % \begin{equation*}
% %     \lim_{n\to \infty}\int_{0}^{\infty} f_n(z) \varphi(z)dz = \int_{0}^{\infty} \lim_{n \to \infty}f_n(z) \varphi(z)dz = \int_{0}^\infty f(z)\varphi(z)dz = 0
% % \end{equation*}
% % which is equivalent to $\lim_{x\to \infty}I(x) = \int_0^{\infty} \lim_{x\to \infty}\alpha(x, x+\epsilon z)\varphi(z) = 0$.
% % \end{proof}


% % \begin{remark}
% % \label{rk: extension remark}
% %     It is trivial to see that $I(x)$ in Lemma \ref{lemma: convergence of acceptance probability at the tail, general light-tailed target} corresponds to \textcircled{1} for $c \to \infty$ in Lemma \ref{lemma: convergence of acceptance probability at the tail, 1D canonical Gaussian}, and both converges to 0 when $x \to \infty$ implies that we can extend our discussion for the 1D canonical Gaussian target to such more general form of light-tailed target densities.
% % \end{remark}

% % \begin{proposition}
% %     With the region $R$ with boundary $K$ defined in \eqref{eqn: region R}, also the target being as the 1D canonical Gaussian, we have $\lim_{x\to \pm \infty}P(|X_t|\leq K \mid X_0=x) = 0$.
% % \end{proposition}

% % \begin{proof}
% % We shall prove this proposition by induction on $t$. When $t=1$, with $X_1 = X_0 + \epsilon z = x+\epsilon z, z \sim N(0,1)$ we write
% % \begin{equation*}
% % \begin{aligned}
% %     P(|X_1| \leq K \mid X_0 = x) &= P(-K \leq x+\epsilon z \leq K)\\
% %     &= \Phi\left(\frac{K-x}{\epsilon}\right) - \Phi\left(\frac{-K-x}{\epsilon}\right) \to 0
% % \end{aligned}
% % \end{equation*}
% % when $x \to \pm \infty$ obviously, with fixed $K>0$ and stepsize $\epsilon > 0$. Now we assume the claim holds at $t$, and define
% % \begin{equation}
% %     f(u) = P(|X_1|\leq K \mid X_0 = u)
% % \end{equation}
% % From the base case we know $f(u) \to 0$ when $u \to \pm \infty$, which equivalently means
% % \begin{equation}
% % \label{eqn: appx proof limit form}
% %     \forall \kappa > 0, \exists \delta = \delta(\kappa) > 0, \forall u \text{ s.t. }|u| > \delta, f(u) < \kappa
% % \end{equation}
% % and now we fix that $\kappa$ (so as related $\delta = \delta(\kappa)$). At $t+1$, we can use $f$ to write
% % \begin{equation}
% % \begin{aligned}
% % P(|X_{t+1}| \leq K\mid X_0=x) &= \mathbb{E}_{X_t}(f(X_t)\mid X_0 = x)\\
% % &= \mathbb{E}(f(X_t)\mathbf{1}(|X_t| > \delta(\kappa))\mid X_0=x) + \mathbb{E}(f(X_t)\mathbf{1}(|X_t| \leq \delta(\kappa))\mid X_0=x)
% % \end{aligned}
% % \end{equation}
% % Note that by \eqref{eqn: appx proof limit form} we see $\mathbb{E}(f(X_t)\mathbf{1}(|X_t| > \delta(\kappa))\mid X_0=x) < \kappa$ for any arbitrary $\kappa>0$ (as here we have $|X_t| > \delta(\kappa) \implies f(X_t) < \kappa$), while
% % \begin{equation*}
% %     \mathbb{E}(f(X_t)\mathbf{1}(|X_t| \leq \delta(\kappa))\mid X_0=x) \leq \left[\sup_{|u| \leq \delta(\kappa)} f(u)\right] P(|X_t|\leq \delta(\kappa) \mid X_0=x) \to 0
% % \end{equation*}
% % based on the inductive hypothesis that $P(|X_t|\leq K \mid X_0=x) \to 0$ for any arbitrary $K > 0$, and here we just set $K=\delta(\kappa)$. Together, we get $P(|X_{t+1}| \leq K\mid X_0=x) \to 0$ when $x\to \pm \infty$, so our claim also holds at $t+1$, and then by induction, the claim is proved.
% % \end{proof}

% % \subsection{Computation of related terms}
% % \label{appx subsec: computation of related terms}
% % This section of the supplementary material is used to give detailed calculations of all necessary terms that are also mentioned in section \ref{subsubsec: first-step RWM and LGWM} and section \ref{subsubsec: multi-step RWM and LGWM}.

% % \paragraph{Explicit expression of \textcircled{1}}
% % \textcircled{1} is defined as, for some $c>x>0$,
% % % \begin{equation}
% % % \label{eqn: expression of textcircled 1}
% % % \begin{aligned}
% % %     &\textcircled{1} = \int_x^{c} \alpha(x, y)q(x,dy) = \int_x^c \exp \left(\frac{1}{2}(x^2 - y^2) \right) \frac{1}{\epsilon\sqrt{2\pi}}\exp\left(-\frac{1}{2\epsilon^2}(y-x)^2\right)dy\\
% % %     &= \frac{1}{\epsilon \sqrt{2\pi}}\int_x^c \exp\left(\frac{1}{2}x^2(1 - \epsilon^{-2})\right) \exp\left(-\frac{1}{2\frac{\epsilon^2}{1 + \epsilon^2}}\left(y - \frac{1}{\epsilon^2(1 + \epsilon^{-2})}x \right)^2\right)\exp\left(\frac{1}{2\epsilon^2(1+\epsilon^2)}x^2\right) dy\\
% % %     &= \frac{1}{\epsilon \sqrt{2\pi}} \exp\left(\frac{\epsilon^2}{2(1+\epsilon^2)}x^2\right) \int_{x - \frac{1}{\epsilon^2(1 + \epsilon^{-2})}x}^{c - \frac{1}{\epsilon^2(1 + \epsilon^{-2})}x}  \exp\left(-\frac{u^2}{2\frac{\epsilon^2}{1 + \epsilon^2}}\right)du \quad \text{by $u = y - \frac{1}{\epsilon^2(1 + \epsilon^{-2})}x$}\\
% % %     &= \frac{1}{\epsilon \sqrt{2\pi}} \exp\left(\frac{\epsilon^2}{2(1+\epsilon^2)}x^2\right) \int_{g_2(x)}^{g_1(c, x)}  \exp\left(-\frac{t^2}{2}\right)\frac{\epsilon}{\sqrt{1+\epsilon^2}}dt \quad \text{by $t = \frac{u}{\epsilon / \sqrt{1 + \epsilon^2}}$}\\
% % %     &= \frac{1}{\sqrt{1 + \epsilon^2}} \exp\left(\frac{\epsilon^2}{2(1+\epsilon^2)}x^2\right) \int_{g_2(x)}^{g_1(c, x)} \frac{1}{\sqrt{2\pi}}\exp\left(-\frac{t^2}{2}\right)dt\\
% % %     &= \frac{1}{\sqrt{1 + \epsilon^2}} \exp\left(\frac{\epsilon^2}{2(1+\epsilon^2)}x^2\right) \left[\Phi(g_1(c, x)) - \Phi(g_2(x))\right]
% % % \end{aligned}
% % % \end{equation}
% % % where $\Phi$ is the CDF for canonical normal distribution $N(0,1)$, and
% % % \begin{align}  
% % % \label{eqn: expression of g1}
% % %     g_1(c, x) &= \frac{\sqrt{1+\epsilon^2}}{\epsilon}\left(c - \frac{1}{\epsilon^2(1 + \epsilon^{-2})}x\right) = \frac{\sqrt{1+\epsilon^2}}{\epsilon}\left(c-\frac{x}{1+\epsilon^2}\right)\\
% % % \label{eqn: expression of g2}
% % %     g_2(x) &= \frac{\sqrt{1+\epsilon^2}}{\epsilon}\left(x - \frac{1}{\epsilon^2(1 + \epsilon^{-2})}x\right) = \frac{\epsilon}{\sqrt{1+\epsilon^2}}x
% % % \end{align}
% % % When $c \to \infty$, we intuitively denote
% % % \begin{equation}
% % %     \lim_{c\to \infty}\textcircled{2} = \frac{1}{\sqrt{1 + \epsilon^2}} \exp\left(\frac{\epsilon^2}{2(1+\epsilon^2)}x^2\right) \left[1 - \Phi(g_2(x))\right]
% % % \end{equation}


% % \paragraph{Limiting behaviour of \textcircled{1}} For \textcircled{1} evaluated in \eqref{eqn: expression of textcircled 1}, when $x \to \infty$ we will have 
% % \begin{equation}
% %     \lim_{x\to \infty}\textcircled{1} = 0
% % \end{equation}
% % this can be deduced by using the standard Gaussian tail asymptotic (Mills ratio) that $1-\Phi(t)\sim \varphi(t)/t$ when $t \to \infty$, so the exponential terms will cancel and remaining a decay term $O(x^{-1})$.

% % \begin{remark}
% %     Note that $\lim_{x\to \infty}\textcircled{1}$ is actually a shorthand notation for $\lim_{x\to \infty} (\lim_{c \to \infty}\textcircled{1})$, as for \textcircled{1} we require $c>x$.
% % \end{remark}

% % \subsection{Discussion for starting position $x \leq 0$}
% % \label{appx subsec: discussion of x<0}
% % In section \ref{subsubsec: first-step RWM and LGWM}, we complete our derivation of statement \eqref{eqn: TV convergence statement, LGWM and RWM} under $x > 0$. To prove this statement at the starting point $x\leq 0$ and when $x \to -\infty$ (so $p = -\mathrm{sgn}(x) = 1$), we proceed identically but swap the sign-test on $z$: Draw $z\sim\mathcal N(0,1)$ and set the proposals $x'= x + \varepsilon z$ for RWM and $w'= x + \epsilon p|z|$ for LGWM. Now, if $z>0$ then $x'=w'$; in this event accept the common proposal for both chains with probability $\alpha(x,x')$, setting $X_1=W_1=x'$ and keeping $P_1=p$ when accepted, otherwise set $X_1=x$ and $(W_1,P_1)=(x,-p)$; on the other hand if $z\leq 0$, the proposals differ and on rejection we keep $X_1=W_1=x$, with $P_1 = -p$. Thus, for $x<0$, the position variables agree after one step either when $z>0$ and the common proposal is accepted, or when both proposals are rejected. To make the laws equal in general, one then shows the acceptance probability vanishes in the tail (now as $x\to-\infty$) for the complementary sign of $z$, which simply follows Lemma \ref{lemma: convergence of acceptance probability at the tail, 1D canonical Gaussian} and the computation of \textcircled{1}, as by symmetry we have
% % \begin{equation*}
% %     \int_x^{\infty} \alpha(x,y)q(x,dy) = \int_{-\infty}^{-x} \alpha(x,y)q(x,dy)
% % \end{equation*}


% % we need to change the coupling algorithm \ref{alg: coupling construction, RWM and LGWM} into the algorithm \ref{alg: coupling construction, RWM and LGWM, x<0}. The deduction would follow similarly as the case $x>0$ described in section \ref{subsubsec: first-step RWM and LGWM}.

% % \begin{algorithm}[ht]
% % \caption{Coupling construction algorithm, RWM and LGWM, for $x\leq 0$}\label{alg: coupling construction, RWM and LGWM, x<0}
% % \begin{algorithmic}[1]

% % \State{Given the starting position $x$, starting momentum $p=-\mathrm{sgn}(x)$.}

% % \State{Sample $z \sim N(0,1)$, compute $x' = x + \epsilon z$ and $w' = x + \epsilon p|z|$.}


% % \If{$z > 0$}
% % \Comment{for this case we have $x'=w'=x+\epsilon |z|$, with $x\leq 0 \implies p=1$}
% % \State{Let $X=W=x'$ with probability $\alpha(x, p, x') = \alpha(x,x')$, otherwise stay.}

% % \Else
% % \State{Let $X=x'$ with probability $\alpha(x,x')$, otherwise stay.}
% % \State{We keep $W=x$ at the original position.}
% % \EndIf

% % \end{algorithmic}
% % \end{algorithm}

% % \subsection{Additional proofs}
% % \label{appx subsec: additional proofs}
% % We will prove $\lim_{x\to \pm \infty}P(|X_t|\leq K \mid X_0=x) = 0$ here by induction on $t$.

% % \begin{proof}
% % Base case: when $t=1$, with $X_1 = X_0 + \epsilon z = x+\epsilon z, z \sim N(0,1)$ we write
% % \begin{equation*}
% % \begin{aligned}
% %     P(|X_1| \leq K \mid X_0 = x) &= P(-K \leq x+\epsilon z \leq K)\\
% %     &= \Phi\left(\frac{K-x}{\epsilon}\right) - \Phi\left(\frac{-K-x}{\epsilon}\right) \to 0
% % \end{aligned}
% % \end{equation*}
% % when $x \to \pm \infty$ obviously, with fixed $K>0$ and stepsize $\epsilon > 0$. Now we assume the claim holds at $t$, and define
% % \begin{equation}
% %     f(u) = P(|X_1|\leq K \mid X_0 = u)
% % \end{equation}
% % From the base case we know $f(u) \to 0$ when $u \to \pm \infty$, which equivalently means
% % \begin{equation}
% % \label{eqn: appx proof limit form}
% %     \forall \kappa > 0, \exists \delta = \delta(\kappa) > 0, \forall u \text{ s.t. }|u| > \delta, f(u) < \kappa
% % \end{equation}
% % and now we fix that $\kappa$ (so as related $\delta = \delta(\kappa)$). At $t+1$, we can use $f$ to write
% % \begin{equation}
% % \begin{aligned}
% % P(|X_{t+1}| \leq K\mid X_0=x) &= \mathbb{E}_{X_t}(f(X_t)\mid X_0 = x)\\
% % &= \mathbb{E}(f(X_t)\mathbf{1}(|X_t| > \delta(\kappa))\mid X_0=x) + \mathbb{E}(f(X_t)\mathbf{1}(|X_t| \leq \delta(\kappa))\mid X_0=x)
% % \end{aligned}
% % \end{equation}
% % Note that by \eqref{eqn: appx proof limit form} we see $\mathbb{E}(f(X_t)\mathbf{1}(|X_t| > \delta(\kappa))\mid X_0=x) < \kappa$ for any arbitrary $\kappa>0$ (as here we have $|X_t| > \delta(\kappa) \implies f(X_t) < \kappa$), while
% % \begin{equation*}
% %     \mathbb{E}(f(X_t)\mathbf{1}(|X_t| \leq \delta(\kappa))\mid X_0=x) \leq \left[\sup_{|u| \leq \delta(\kappa)} f(u)\right] P(|X_t|\leq \delta(\kappa) \mid X_0=x) \to 0
% % \end{equation*}
% % based on the inductive hypothesis that $P(|X_t|\leq K \mid X_0=x) \to 0$ for any arbitrary $K > 0$, and here we just set $K=\delta(\kappa)$. Together, we get $P(|X_{t+1}| \leq K\mid X_0=x) \to 0$ when $x\to \pm \infty$, so our claim also holds at $t+1$, and then by induction, the claim is proved.
% % \end{proof}
%\input{Notes/appendix_full}

\end{document}
